\documentclass[output=paper,colorlinks,citecolor=brown]{langscibook}
\ChapterDOI{10.5281/zenodo.20611903}
\author{Jerzy Gaszewski\orcid{}\affiliation{University of Lodz}}

\title{Uninflectedness as a rule in Polish, an inflected language}

\abstract{The chapter describes uninflected nouns referring to women in Polish. The absence of inflection on these nouns results from the lack of an inflectional paradigm that would suit the class in terms of both morphonology and semantics. In this case, uninflectedness, generally very peculiar in the language, has developed into a signal of feminine gender in its own right and given rise to a productive grammatical pattern. The lack of marking allows distinguishing the feminine words and the normally inflecting masculine counterparts. The text presents the functioning of the class that so poorly fits the Polish system, heavily reliant on inflection, competing patterns that do not involve uninflectedness as well as constraints on the use of the uninflected words. The last section gives a brief historical overview of the emergence of the category and shows that the process happened twice: first for surnames, then for personal common nouns. The article also addresses the ongoing debate among native speakers and recent developments which could potentially trigger the downfall of the discussed uninflected pattern.}


\IfFileExists{../localcommands.tex}{
   \addbibresource{../localbibliography.bib}
   \usepackage{tabularx,multicol}
\usepackage[hyphens]{url}
\urlstyle{same}

\usepackage{langsci-optional}
\usepackage{langsci-lgr}
\usepackage{langsci-gb4e}

\usepackage{soul}
% \usepackage{booktabs}
% \usepackage{geometry}
\usepackage{multirow}
% \usepackage{makecell} %safely breaks a line inside a table cell
% \usepackage{tablefootnote} %footnotes in table captions
\usepackage{enumerate}
\usepackage{tikz}
\usepackage{array}
\usepackage[shortlabels]{enumitem}

%for having several figures on one line, turning words in 90°angles
% \usepackage{graphicx}
\usepackage{subcaption}

% \usepackage{caption} %for making the caption line longer in concrete cases

% \usepackage{rotating} %rotate a whole page; for large tables

\usepackage{fontspec}
\newfontfamily\charis{Charis SIL} %for being able to use the symbol ‿ and a better command of some of the boldfaced diacritics in 01
\newfontfamily\japanesefont{Noto Serif CJK JP} % for Japanese script
% \usepackage{tipa}

% \usepackage{needspace} %to keep exs on one page

\usepackage{xcolor}

\usepackage{pgfplots}
\pgfplotsset{compat=1.18}

%%%%%   AVMs for 02	%%%%%
\usepackage{langsci-avm}
\avmsetup{switch = ?}


\usepackage{langsci-branding}

   %\newcommand*{\orcid}{}


\makeatletter
\let\thetitle\@title
\let\theauthor\@author
\makeatother

\newcommand{\togglepaper}[1][0]{
   \bibliography{../localbibliography}
   \papernote{\scriptsize\normalfont
     \theauthor.
     \titleTemp.
     To appear in:
     E. Di Tor \& Herr Rausgeberin (ed.).
     Booktitle in localcommands.tex.
     Berlin: Language Science Press. [preliminary page numbering]
   }
   \pagenumbering{roman}
   \setcounter{chapter}{#1}
   \addtocounter{chapter}{-1}
}

\newbool{bookcompile}
\booltrue{bookcompile}
\newcommand{\bookorchapter}[2]{\ifbool{bookcompile}{#1}{#2}}



\renewcommand{\thempfootnote}{\fnsymbol{mpfootnote}} % use * for footnotes in float env

% \newcolumntype{L}[1]{>{\raggedright\arraybackslash}p{#1}} %left-aligned (non-justified) text in tables
% \newcolumntype{R}[1]{>{\raggedleft\arraybackslash}p{#1}} % right-aligned (non-justified) text in tables

\definecolor{customgreen}{RGB}{27,158,119}
\definecolor{custompink}{RGB}{231,42,138}
\definecolor{customblue}{RGB}{117,112,179}
\definecolor{customlightblue}{RGB}{143,170,220}
\definecolor{customdarkblue}{RGB}{68,114,196}
\definecolor{customyellow}{RGB}{225,192,0}
\definecolor{customorange}{RGB}{233,137,21}
\definecolor{customgray}{RGB}{229,229,229}


% Left-aligned indexes
% % \makeatletter
% % \patchcmd{\theindex}{\twocolumn}{\raggedright\twocolumn}{}{}
% % \makeatother

%to make Japanese scripts in the bibliography smaller
\newcommand{\japbib}[1]{{\japanesefont\small #1}}

%%%%%%%%%%%%%%%%%%%	MACROS for 02	%%%%%%%%%%%%%%%%

%%% Formatting	%%%%%

\newcommand{\textex}[1]{\textit{#1}} 		 %for formatting linguistic examples in-text%
%command-shift X

\newcommand{\lxm}[1]{\textsc{#1}}  		%for formatting names of lexemes
%command-option-shift L

%%%%%%%%% Abbreviations	%%%%%%

\newcommand{\featname}[1]{\textsc{#1}}
\newcommand{\featspec}[2]{\featname{#1}:{#2}}
\newcommand{\feat}[2]{\textsc{#1}:{#2}}  				% \feat{FEATURE}{value}  sc's
\newcommand{\featnameonly}[1]{\textsc{#1}} 			% \featnameonly{featurename}  sc's

\newcommand{\pr}{$^{\prime}$}

\newcommand{\Corr}{\textit{Corr}}
\newcommand{\propm}{\textit{pm}}

\newcommand{\PC}{Π$_{C}$}

\newcommand{\PF}{Π$_{F}$}
\newcommand{\PR}{Π$_{R}$}

\newcommand{\lex}{$\mathcal{L}$}

   %% hyphenation points for line breaks
%% Normally, automatic hyphenation in LaTeX is very good
%% If a word is mis-hyphenated, add it to this file
%%
%% add information to TeX file before \begin{document} with:
%% %% hyphenation points for line breaks
%% Normally, automatic hyphenation in LaTeX is very good
%% If a word is mis-hyphenated, add it to this file
%%
%% add information to TeX file before \begin{document} with:
%% %% hyphenation points for line breaks
%% Normally, automatic hyphenation in LaTeX is very good
%% If a word is mis-hyphenated, add it to this file
%%
%% add information to TeX file before \begin{document} with:
%% \include{localhyphenation}
\hyphenation{
    Al-ex-an-der
    Ber-ber
    chan-ges
    cir-cum-stan-ces
    coin-age
    Fraj-zyngier
    ita-liano
    Ja-pa-nese
    Kö-nig
    Krzy-żanowski
    li-te-ra-tur-no-go
    mi-zen-kei
    par-a-digm
    phe-nom-e-non
    ren-tai-shi
    Slo-vene
    Spen-cer
    Stoj-kov
    Ver-bi-ca-re-se
    wide-spread
}

\hyphenation{
    Al-ex-an-der
    Ber-ber
    chan-ges
    cir-cum-stan-ces
    coin-age
    Fraj-zyngier
    ita-liano
    Ja-pa-nese
    Kö-nig
    Krzy-żanowski
    li-te-ra-tur-no-go
    mi-zen-kei
    par-a-digm
    phe-nom-e-non
    ren-tai-shi
    Slo-vene
    Spen-cer
    Stoj-kov
    Ver-bi-ca-re-se
    wide-spread
}

\hyphenation{
    Al-ex-an-der
    Ber-ber
    chan-ges
    cir-cum-stan-ces
    coin-age
    Fraj-zyngier
    ita-liano
    Ja-pa-nese
    Kö-nig
    Krzy-żanowski
    li-te-ra-tur-no-go
    mi-zen-kei
    par-a-digm
    phe-nom-e-non
    ren-tai-shi
    Slo-vene
    Spen-cer
    Stoj-kov
    Ver-bi-ca-re-se
    wide-spread
}

   \boolfalse{bookcompile}
   \togglepaper[23]%%chapternumber
}{}

\begin{document}
\maketitle

\section{Introduction} 
\il{Polish|(}\is{masculine|(}

Polish is a language with rich \isi{inflection}, involving nouns, pronouns, determiners, adjectives, participles, numerals and verbs. Yet, there are nouns (1\% of all, cf.~\citealt[159]{stefanczyk2007}) where the expected marking is lacking. This paper focuses on one type of uninflected\is{uninflectedness} nouns in Polish, originally masculine words used to refer to women (called \textsc{U-forms}\is{U-form} here, cf. Section \ref{11sec3.1}). Their \isi{uninflectedness}, otherwise exceptional in the language, is essentially regular. Most of the words in question are borrowings, but we also find native ones in this group. Let us add that the said nouns are originally regularly inflecting\is{inflection} words referring to men and having the masculine grammatical \isi{gender}.\footnote{Polish is said to have five \isi{gender} categories: masculine personal (virile\is{masculine!virile}), masculine animate/animal, masculine inanimate, \isi{feminine} and \isi{neuter}. The plain label ``masculine" will still appear in the paper. With words referring to people it means masculine personal. The three masculine genders\is{gender} are also syncretised in a number of grammatical positions. Let us stress that there is no grammatical context in Polish that would distinguish as many as five \isi{gender} values. The five-way division proposed by \citet{manczak1956} results from the superposition of the maximal distinctions made	in the \isi{accusative} singular and plural\is{masculine|)}.\label{11ftn1}}

The paper describes how this class of words functions in a language whose grammatical system it seems to fit so poorly, esp. considering the existence of rival morphological patterns. We will also sketch the historical development of U-forms\is{U-form}, an open category at odds with basic Polish grammar, a regular pattern that as a whole is an exception.

Some aspects of the topic are socially very salient and constitute the \isi{subject} of an ongoing public debate among native speakers. However, the bone of contention is not the \isi{uninflectedness} of U-forms\is{U-form}. It is the rival inflectable\is{inflection} forms referring to women (which we call \textsc{D-forms}\is{D-form}) whose increasingly common use (for the very same concepts) is controversial. U-forms\is{U-form} surface in the debate, but as a peripheral issue. The roles are reversed here; our focus is on U-forms\is{U-form} as a particularly interesting \isi{case} of \isi{uninflectedness} and the hotly debated D-forms\is{D-form} are merely a point of reference. Thus, the paper is not meant to be a voice in the current debate in Poland; the aim is to analyse the usage without taking sides. We refer to the debate because it is a relevant part of the sociolinguistic reality.

{\sloppy Many examples come from the National Corpus of Polish (NKJP, \citealt{przepiorkowskibankogorskilewandowskatomaszczyk2012}), but the paper lacks an overall statistical analysis of the phenomena described. This is due to space limitations and particular difficulties with such a task. The tagging in the NKJP (or any corpus of Polish) logically relies on inflectional categories\is{inflection}, which does not allow extracting reliable results for whole classes of uninflected nouns\is{uninflectedness}. However, the expressions referring to frequency used here are not mere opinions or guesses. First of all, they reflect the relative difficulty in finding examples. Furthermore, to compensate for the lack of the overall statistics, numbers for selected example phrases are included to support the points made at several points in the discussion. The related queries in the corpus were performed with the use of the PELCRA search engine \citep{pezik2012}.\par}

Section \ref{11sec2} of the paper briefly presents the Polish noun declension, then moves to factors conducive to \isi{uninflectedness} of nouns. Section \ref{11sec3} is an analysis of the functioning of the uninflected words\is{uninflectedness} for women in the language at present (21\textsuperscript{st}~century). Section \ref{11sec4} provides a historical overview of the development of uninflected\is{uninflectedness} and inflectable\is{inflection} words for women, including current changes.

\section{Declension and lack thereof} \label{11sec2}

\subsection{Declension in Polish} \label{11sec2.1}

This section is meant to present Polish declension without delving into detail. For a comprehensive picture the readers are referred to grammars \citep{zagorska-brooks1975, laskowski1979, bielec1998, swan2002}. Let us start with an example of declension in use. (\ref{11ex1}) is the opening sentence of the preface to one academic grammar of Polish, for clarity divided into clauses here.

\begin{exe}
	\ex \label{11ex1}
	\begin{xlist}
		\ex \label{11ex1a}
		\gll Książk-a t-a powstała na zamówieni-e. \\
		book-\textsc{nom.sg} this-\textsc{nom.sg.f} arise.\textsc{pst.f} \textsc{prep} commission-\textsc{acc.sg} \\
		\trans ‘This book was commissioned.’
		
		\ex \label{11ex1b}
		\gll Nie odmawia się Wydawnictw-u, \\
		\textsc{neg} refuse.\textsc{3sg} \textsc{impers} publishing\_house-\textsc{dat.sg} \\
		\trans ‘One does not refuse a publishing house,’
		
		\ex \label{11ex1c}
		\gll w któr-ym druk-Ø jest dla autor-a honor-em. \\
		in which-\textsc{loc.sg.n} print-\textsc{nom.sg} be.\textsc{3sg} for author-\textsc{gen.sg} honour-\textsc{ins.sg} \\
		\trans ‘to publish with which it is an honour for the author.’
		
		\hfill (\citealt[9]{nagorko1997})
	\end{xlist}
\end{exe}

(\ref{11ex1}) features a fairly representative selection of noun forms. We have a typical \isi{masculine} \isi{nominative} with a zero suffix, \emph{druk} in (\ref{11ex1c}), a typical \isi{feminine} one with the suffix \emph{-a}, \emph{książka} in (\ref{11ex1a}), a neuter \isi{accusative} identical in form with the \isi{nominative}, \emph{zamówienie} in (\ref{11ex1a}), with a non-zero marker \emph{-e}. There are distinct oblique \isi{case} forms: a \isi{dative}, \emph{wydawnictwu} in (\ref{11ex1b}) and an \isi{instrumental}, \emph{honorem} in (\ref{11ex1c}), as well as a \isi{masculine} personal \isi{genitive} identical to the \isi{accusative}, \emph{autora} in (\ref{11ex1c}). The \isi{locative} is exemplified by the relative pronoun \emph{którym} in (\ref{11ex1c}). In sum, all Polish cases\is{case} appear in (\ref{11ex1}), except the \isi{vocative}, which is unlikely in academic prose. However, at the beginning of a preface one could easily imagine one of the \isi{vocative} phrases in (\ref{11ex2}) preceding (\ref{11ex1}).

\begin{exe}
	\ex \label{11ex2}
	\begin{xlist}
		\ex \label{11ex2a}
		\gll drog-i Czytelnik-u \\
		dear-\textsc{voc.sg.m} reader-\textsc{voc.sg} \\
		\trans ‘dear Reader’
		
		\ex \label{11ex2b}
		\gll drodz-y Czytelnic-y \\
		dear-\textsc{voc.pl.vir} reader-\textsc{voc.pl} \\
		\trans ‘dear Readers’
		
		\ex \label{11ex2c}
		\gll drogi-e Czytelnicz-k-i, drodz-y Czytelnic-y \\
		dear-\textsc{voc.pl.nvir} reader-\textsc{f-voc.pl} dear-\textsc{voc.pl.vir} reader-\textsc{voc.pl} \\
		\trans ‘dear Readers [female], dear Readers [male]’
	\end{xlist}
\end{exe}

As for syntactic relations, the \isi{nominative} forms in (\ref{11ex1}) are subjects\is{subject} of the clauses, the \isi{instrumental} marks the \isi{subject} complement after a copula, the \isi{accusative} is governed by the preposition \emph{na}, etc. Case\is{case} marking fulfils the grammatical requirements of the words the nouns are structurally connected with. This phenomenon, syntactic \isi{government} (in the traditional sense, cf. \citealt[246--250]{matthews1981}), relies on the category of \isi{case} and is essential for signalling clause structure. Case\is{case} (as well as \isi{number} and \isi{gender}) is also marked on words modifying nouns, as in the demonstrative \emph{ta} in (\ref{11ex1a}) and the adjectives in (\ref{11ex2}). This kind of \isi{agreement} is crucial for the cohesion of noun phrases. Agreement\is{agreement} can also be seen on the verb \emph{powstała} in (\ref{11ex1a}), which shows the \isi{number} as well as \isi{gender} of the \isi{subject} of the clause. Gender\is{gender} is only marked on verb forms based on the past tense, historically a participle.

Polish nouns inflect\is{inflection} for 7 cases\is{case} and 2 numbers\is{number}. This yields a 14-slot system, but because of \isi{case} syncretism, mostly systematic, as in (\ref{11ex1}--\ref{11ex2}), no noun has more than eleven distinct forms and many have fewer. Adjectives, participles and determiners show the \isi{case} and \isi{number} of the noun they modify and additionally mark its \isi{gender}, an inherent property of a given noun. The declension markers are solely suffixes, although \isi{inflection} may be accompanied by sound alternations in the stem, most of them predictable, cf. the forms in (\ref{11ex2}). The suffixes are fusional; it is never possible to separate the marking of \isi{case} and \isi{number} (and \isi{gender} in noun modifiers). Nouns fall into a \isi{number} of declension classes which have different sets of markers for specific grammatical values. If this general description is reminiscent of classical Indo-European languages such as \ili{Latin} or \ili{Sanskrit}, the association is correct -- the overall makeup of the Polish declension system is conservative.

Which declension class a given noun belongs to is the result of the interplay of its morphonetic shape and grammatical \isi{gender}, though the two factors are correlated to a considerable extent. The \isi{nominative} singular of a noun (its citation form) usually reveals the \isi{gender} in the traditional three-way division. Most \isi{masculine} nouns end in a consonant (the suffix being \emph{-Ø}), as in (\ref{11ex1c}), most feminines\is{feminine} end in \emph{-a}, as in (\ref{11ex1a}), and most neuters\is{neuter} end in \emph{-o} or \emph{-e}, as in (\ref{11ex1a}). This correlation is not absolute, but it delineates the major declension paradigms\is{paradigm}.\footnote{They are ``major" in terms of the \isi{number} of nouns, but also the ease of acquisition. \citet[625--626]{smoczynska1985} shows that nouns of minor classes are problematic for young speakers.} The three \isi{masculine} genders\is{gender} are normally distinguished by different patterns of syncretism in the \isi{accusative} forms (matched by different \isi{agreement} markers, cf. fn. \ref{11ftn1}).

The defining parameter for \isi{gender} is the distinctions of \isi{agreement} markers on noun modifiers (\citealt[231]{hockett1958} after \citealt[1]{corbett1991}, cf. also fn. \ref{11ftn1}). Thus, a Polish noun has a set value of this category, which manifests itself in the \isi{inflection} of noun modifiers and verb forms based on the past tense. As previously mentioned, in practice \isi{gender} is usually deducible from the form of the noun itself. The labels ``\isi{masculine}" and ``\isi{feminine}" imply a semantic basis for the \isi{gender} system, not really strong in Polish \citep[175]{kucala1978}, but still perceivable. Gender\is{gender} is mostly a formal grammatical phenomenon apart from the opposition of \isi{masculine} personal and \isi{feminine} in nouns for people, which will prove crucial here.

To further complicate the assignment of declension class, the morphonetic shape concerns not only the very final sound of the basic form, but also the last sound of the stem (which is word-final if the suffix is zero). Stem-final consonants may be \textit{hard}, i.e. non-palatalised, or \textit{soft}, i.e. palatal(ised), there are also intermediate categories. Suffixes for some grammatical values depend on this division. Nouns with different types of stems are distributed across all genders\is{gender}. Table \ref{11tab1} exemplifies the major \isi{feminine} paradigms\is{paradigm} and the hard-stem \isi{masculine} personal \isi{paradigm} for comparison.

\begin{table}[ht]
	\centering
	\begin{tabularx}{\textwidth}{llllll}
		\lsptoprule
			& mama & ciocia & córka & uczennica & doktor \\
			& ‘mum’ & ‘aunt’ & ‘daughter’ & ‘pupil, \textsc{f}’ & ‘doctor, \textsc{m}’ \\
			\midrule
			& \textsc{f} hard & \textsc{f} soft & \textsc{f} velar & \textsc{f} historically & \textsc{vir} hard \\
			& stem & stem & stem & soft stem & stem \\
			\midrule
			\textsc{nom sg} & mam-a & cioci-a & córk-a & uczennic-a & doktor-Ø \\
			\textsc{gen sg} & mam-y & cioc-i & córk-i & uczennic-y & doktor-a \\
			\textsc{dat sg} & mami-e & cioc-i & córc-e & uczennic-y & doktor-owi \\
			\textsc{acc sg} & mam-ę & cioci-ę & córk-ę & uczennic-ę & doktor-a \\
			\textsc{ins sg} & mam-ą & cioci-ą & córk-ą & uczennic-ą & doktor-em \\
			\textsc{loc sg} & mami-e & cioc-i & córc-e & uczennic-y & doktorz-e \\
			\textsc{voc sg} & mam-o & cioci-u & córk-o & uczennic-o & doktorz-e \\
			\midrule
			\textsc{nom pl} & mam-y & cioci-e & córk-i & uczennic-e & doktorz-y \\
			\textsc{gen pl} & mam-Ø & cioć-Ø & córek-Ø & uczennic-Ø & doktor-ów \\
			\textsc{dat pl} & mam-om & cioci-om & córk-om & uczennic-om & doktor-om \\
			\textsc{acc pl} & mam-y & cioci-e & córk-i & uczennic-e & doktor-ów \\
			\textsc{ins pl} & mam-ami & cioci-ami & córk-ami & uczennic-ami & doktor-ami \\
			\textsc{loc pl} & mam-ach & cioci-ach & córk-ach & uczennic-ach & doktor-ach \\
			\textsc{voc pl} & mam-y & cioci-e & córk-i & uczennic-e & doktorz-y \\
		\lspbottomrule
	\end{tabularx}
	\caption{Several major declension paradigms\is{paradigm}}
	\label{11tab1}
\end{table}

Several minor paradigms\is{paradigm} include nouns falling outside of the scope of the correlation described above. There may be a ``mismatch'' of the \isi{nominative} singular suffix and \isi{gender}, as in \isi{masculine} nouns ending in \emph{-a}, e.g. \emph{władca} `ruler' and \emph{mężczyzna} `man' or \isi{feminine} nouns ending in a consonant, e.g, \emph{mysz} `mouse' and \emph{ciemność} `darkness'. The \isi{nominative} singular may end in a way not covered by the general rule, as in \isi{feminine} nouns in \emph{-i}, e.g. the honorific \emph{pani} `madam, lady, Mrs' and \emph{władczyni} `female ruler'. Still other nouns use the distinct suffixes of adjectives, e.g. \emph{szeregowy} `private (military rank), \textsc{vir}', \emph{królowa} `queen, \textsc{f}', \emph{znaleźne} `finder's fee, \textsc{n}'. In general, minor declension classes have a clearer semantic profile than the major ones, cf. the discussion of (\ref{11ex4}) and fn. \ref{11ftn10}. Let us note that many of them are productive patterns due to association with specific derivational suffixes e.g. the agentive \emph{-c(a)} in \emph{władca}, feminative \emph{-yn(i)}/\emph{-in(i)} in \emph{władczyni} or abstract \emph{-ość} in \emph{ciemność} above.

In sum, the main factors influencing noun declension in Polish are \isi{gender} and morphonetics. Some (groups of) words show further idiosyncrasies. The resultant \isi{number} of paradigms\is{paradigm} (exact sets of suffixes in the 14-slot grid) is considerable. However, there are numerous similarities across the paradigms\is{paradigm}. For example, the \isi{feminine} classes in Table \ref{11tab1} share some suffixes in the singular. The suffixes for other cases\is{case} differ, but all four patterns syncretise the \isi{dative} and \isi{locative}. In those represented by \emph{ciocia} and \emph{uczennica}, the \isi{genitive} is also identical to these two cases\is{case}. In the plural, the similarities are even more evident and some suffixes are shared with the \isi{masculine} class (and others). Thus, the numerous exact sets of suffixes for the 14 slots are to a large extent variations on relatively few basic patterns.

\subsection{Lack of declension in Polish} \label{11sec2.2}

Let us now review the factors that make some Polish nouns resist inflection\is{uninflectedness}. Most of them are borrowings, but as we shall see, native words (and new coinages) can also be uninflected\is{uninflectedness}. The first factor is the incompatibility of morphonetic shape of the given word and available declension paradigms\is{paradigm}. Consider words ending in \emph{-u}, which cannot be a \isi{nominative} suffix, so the whole word is interpreted as the stem. This does not help much because such stems (with putative zero suffixes)\footnote{Some grammars of Polish posit zero suffixes for all \isi{paradigm} slots of uninflectable nouns\is{uninflectedness} (e.g.~\citealt[271]{orzechowska1998}). Here, subtler distinctions are necessary. In glosses for normally uninflectable nouns\is{uninflectedness} inflectional categories\is{inflection} (retrievable from the context) are shown in brackets, cf.~(\ref{11ex3}). In the \isi{case} of normally inflectable nouns\is{inflection}, zero suffixes (in opposition to non-zero suffixes on the same words) are used, cf.~(\ref{11ex5}). In clear cases\is{case} of contextual \isi{uninflectedness} no morpheme segmentation is used and the category values given match the form of the noun and not the context, cf.~(\ref{11ex7a}).} still cannot be associated with any pattern of declension. Consequently, loanwords like \emph{malibu} in (\ref{11ex3}), \emph{kakadu} `cockatoo', \emph{Peru}, \emph{Nehru} or \emph{Katmandu} remain uninflectable\is{uninflectedness}.

\begin{exe}
	\ex \label{11ex3}
	\gll butelk-a malibu \\
	bottle-\textsc{nom.sg} malibu[\textsc{gen.sg}] \\
	\trans ‘a bottle of malibu’
\end{exe}

Apart from the final \emph{-u}, stressed final vowels are similarly problematic for the Polish system. These are found in loanwords that preserve the original oxytone pattern and in many acronyms (including native ones), e.g. \emph{menu} [mɛˈni], \emph{Delacroix}, \emph{Calais}, \emph{BBC} or \emph{PKO} (name of a bank).

To make sense of (\ref{11ex4}), we must factor in semantics. Even though nouns with final consonants are commonplace in Polish and they normally inflect\is{inflection}, they are \isi{masculine} (with the suffix \emph{-Ø}) and the corresponding \isi{paradigm} is \isi{masculine}, cf.~(\ref{11ex5}).

\begin{exe}
	\ex \label{11ex4}
	\gll Czy nie przestaje Pan-i być rozsądn-ą bizneswoman? \\
	\textsc{q} \textsc{neg} stop.\textsc{3sg} Mrs-\textsc{nom.sg} be.\textsc{inf} rational-\textsc{ins.sg.f} businesswoman[\textsc{ins.sg}] \\
	\trans ‘Don’t you stop being a rational businesswoman?’
	
	\hfill (NKJP, \citealt{przepiorkowskibankogorskilewandowskatomaszczyk2012})
	
	\ex \label{11ex5}
	\gll Czy nie przestaje Pan-Ø być rozsądn-ym biznesmen-em? \\
	\textsc{q} \textsc{neg} stop.\textsc{3sg} Mr-\textsc{nom.sg} be.\textsc{inf} rational-\textsc{ins.sg.m} businessman-\textsc{ins.sg} \\
	\trans ‘Don’t you stop being a rational businessman?’
\end{exe}

Because of its meaning,\emph{ bizneswoman} in (\ref{11ex4}) must be \isi{feminine} and the \isi{gender} assignment does not allow the noun to decline according to the \isi{masculine} \isi{paradigm}. This is a clear \isi{case} of incompatibility: a potentially matching \isi{paradigm} exists, but it is not applied because of the given noun's semantics, reflected in \isi{gender} assignment.

A similar issue arises with words that have a final \emph{-y} [ɨ], which is a \isi{nominative} singular suffix in the \isi{masculine} \isi{paradigm} for adjectives. Nouns that inflect\is{inflection} this way are typically \isi{masculine} personal, cf. \emph{szeregowy} `private' in Section \ref{11sec2.1}. Foreign male names that match the pattern, e.g. \emph{Harry}, inflect\is{inflection}, while semantically unfit borrowings, e.g. \emph{brandy}, do not.

Neuter\is{neuter} nouns in \emph{-um}, e.g. \emph{kuriozum} `oddity' are a unique group in that they do not inflect\is{uninflectedness} in the singular, but do so in the plural\is{inflection}. The pattern has arisen with \ili{Latin} borrowings, but is productive, as many other minor ones. A very recent (and likely short-lived) example is \emph{miraculum}, a name of type of magical object in \emph{Miraculous: Tales of Ladybug \& Cat Noir}, an animated television series from the 2010s.

Let us stress that when nouns are borrowed into Polish, the original form is taken to be the \isi{nominative} singular. Numerous noun forms have a final \emph{-u}, but it serves as a suffix of oblique cases\is{case}. Borrowed foreign forms like \emph{malibu} are never interpreted as having such a suffix.\footnote{Following this line of thought, (\ref{11ex3}) could be interpreted as having a \isi{genitive} suffix \emph{-u}, found in \isi{masculine} inanimate paradigms\is{paradigm}, yielding a \isi{nominative} form *\emph{malib}. In reality, borrowings are never treated this way and the very idea would seem absurd to native speakers. This fantasy scenario is inspired by the treatment of loanwords in \ili{Kambaata} (Cushitic, Ethiopia), which are always inflected\is{inflection} no matter what phonetic shape they have \citep{treis2023}.} The \isi{nominative} singular functions as the basic form in the Polish grammatical tradition and is used as the entry form in dictionaries. The behaviour of borrowings shows that this convention is supported by the feeling of native speakers.

The second factor promoting lack of inflection\is{uninflectedness} can be called the unfamiliar character of the noun. Even though the place names \emph{Vancouver} and \emph{Hanower} `Hannover' end in exactly the same sounds in Polish, only the latter regularly inflects\is{inflection}. This factor is very elusive and cannot be boiled down to geographical distance or low frequency.\footnote{In the \isi{case} of the two city names, this is completely mistaken. \emph{Vancouver} is actually more frequent in the NKJP than \emph{Hanower}, 4,503 vs. 3,331 attestations, respectively.} Names of many cities in Germany, a neighbour of Poland, do not inflect\is{uninflectedness}, e.g. \emph{Essen} or \emph{Monachium} `Munich'. On the other hand, names of numerous very distant and clearly exotic places do inflect\is{inflection}, e.g. \emph{Osaka} or \emph{Bangkok}.

Similar observations can be made about common nouns for concepts of foreign origin. Many such words are normally uninflectable\is{uninflectedness}, even though possibly matching paradigms\is{paradigm} exist, e.g. \emph{judo}, \emph{karate, reggae, kakao} `cocoa', and \emph{mango}. Other established borrowings do inflect\is{inflection}, however. The following examples are semantically similar to the previous set: \emph{joga} `yoga', \emph{rap}, \emph{cha-cha}, \emph{kawa} `coffee' and \emph{papaja} `papaya'.

To sum up, some foreign words function as quotes from other languages \citep[36, 40]{krzyzanowski2013} and do not inflect\is{uninflectedness}, although which ones behave so is really a matter of language use, which is not uniform either (see below). According to \citet[84--85]{dunaj2004}, the features promoting \isi{uninflectedness} are low frequency of the potential \isi{paradigm} and the given borrowing as such, as well as the learned character of the
word.

We mentioned in passing the idea of using words as quotes. In fact, this is a typical environment for contextual suspension of inflection. Thus, in (\ref{11ex6}) the normally inflectable\is{inflection} word \emph{kwiat} does not adopt the form expected after the preposition \emph{do}. The basic form is used for metalinguistic quoting.

\begin{exe}
	\ex \label{11ex6}
	\gll rym-Ø do ``kwiat'' \\
	rhyme-\textsc{nom.sg} \textsc{prep} flower.\textsc{nom.sg} \\
	\trans ‘a rhyme for [the word] flower’ \hfill (NKJP)
\end{exe}

Note that an inflected\is{inflection} word form can be quoted as well, but then its overt categories are not tied to the actual context. However, the unchanged basic form of declinable words in Polish appears often in structures that are not straightforward quoting, namely \textsc{appositions}\is{apposition}, cf. (\ref{11ex7}).

\begin{exe}
	\ex \label{11ex7}
	\begin{xlist}
		\ex \label{11ex7a}
		\gll nad rzek-ą San \\
		upon river-\textsc{ins.sg} San.\textsc{nom.sg} \\
		
		\ex \label{11ex7b}
		\gll nad San-em \\
		upon San-\textsc{ins.sg} \\
		\trans ‘on the river San’
	\end{xlist}
\end{exe}

The name in (\ref{11ex7a}) remains in the \isi{nominative}, quoted, as it were, and the syntactic functions of the whole phrase are marked on the common noun \emph{rzeka} `river'. This happens despite the name being otherwise inflectable\is{inflection} as in (\ref{11ex7b}). Appositional\is{apposition} phrases where one noun (usually more general and coming first) is inflected\is{inflection} while the other is given in the basic form, even though it is generally inflectable\is{inflection}, constitute a common pattern in modern Polish, which easily incorporates uninflectable\is{uninflectedness} nouns. The phrase \emph{strusi emu} in (\ref{11ex8}) follows the pattern of \isi{apposition} observed in (\ref{11ex7a}). The more general noun is declined and thus shows the grammatical relations of the whole phrase to other words whereas the more specific noun bears no inflectional\is{uninflectedness} marking. The difference is that \emph{emu}, like \emph{malibu} in (\ref{11ex3}), is always uninflectable\is{uninflectedness}, unlike the river name in (\ref{11ex7}).

\begin{exe}
	\ex \label{11ex8}
	\gll przodk-owie dzisiejsz-ych strus-i emu \\
	ancestor-\textsc{nom.pl} today.\textsc{adj-gen.pl} ostrich-\textsc{gen.pl} emu[\textsc{gen.pl}] \\
	\trans ‘the ancestors of modern emus’ \hfill (NKJP)
\end{exe}

For many of the nouns mentioned so far there is variation in usage. This is certainly true about the very fluid concept of unfamiliar character. If a given noun is normally not declined, but its morphonetic shape does not preclude assigning a declension \isi{paradigm}, it is likely that some speakers will inflect it\is{inflection}. For example, \emph{Vancouver} may be uninflectable\is{uninflectedness} in standard usage, but some speakers do decline the name.\footnote{There are 2 attestations of \emph{do Vancouveru} `to Vancouver' in the NKJP and 132 of the standard phrase without inflection\is{uninflectedness}. There are 30 attestations of \emph{w Vancouverze} `in Vancouver' and 1210 of the version without inflection\is{uninflectedness}.}

A historical parallel is the usage of the noun \emph{radio}. Initially, the word was uninflectable\is{uninflectedness}, presumably due to novelty of the invention, coupled with unfamiliar phonetics, the palatalised plosive [d\textsuperscript{j}] being rare in Polish. Today, however, the inflected\is{inflection} forms have become dominant. Thus, some tendency to eliminate particular cases\is{case} of \isi{uninflectedness} can be noted. This process leads to variation in contemporary usage and changes of the properties of lexemes\is{lexeme} over longer stretches of time.

Changes such as that of the noun \emph{radio} are rather expected; if nouns generally decline, it is natural that new ``unfamiliar'' nouns eventually comply with the regular pattern. However, \citet{kurek2019} documents an opposite tendency to avoid declension in present day Polish. One context in which normally inflectable\is{inflection} nouns are used in their basic forms is official documents. This is linked to generating them on the basis of electronic databases. Various data, particularly surnames, are stored in \isi{nominative} forms and then inserted directly into documents, whatever the actual linguistic context \citep[57]{kurek2019}. This kind of ``officialese'' may serve as a model for some speakers in other situations. Uncertainty about the actual forms of particular words also plays a role \citep[61--62]{kurek2019}. Faced with the risk of creating a wrong inflected\is{inflection} form, speakers tend to prefer \isi{uninflectedness} as the safest option, even though this may sometimes obscure the meaning or lead to misunderstandings \citep[157--159]{kurek2019}. Such usage has not been accepted as standard as yet, but is clearly present. The process is most advanced with surnames. These normally appear in \isi{apposition} with given names, which resist the tendency much more and usually inflect\is{inflection}.

\section{The uninflected\is{uninflectedness} words for women in present-day Polish} \label{11sec3}

\subsection{The array of constructions} \label{11sec3.1}

The uninflected\is{uninflectedness} nouns referring to women represent a regular pattern, but are clearly a minority among \isi{feminine} personal nouns in Polish. The examples below present both the uninflected\is{uninflectedness} forms and competing patterns. We shall start by defining the relevant word types. The focal group will be called \textsc{U-forms}\is{U-form} (for `uninflected'\is{uninflectedness}), nouns referring to women that have no inflectional marking\is{uninflectedness} in syntactic positions where non-zero suffixes are expected. The major competing pattern are \textsc{D-forms}\is{D-form} (for `derivation'), nouns referring to women that have clear morphologically related \isi{masculine} counterparts (and so can be construed as derived from them), but inflect\is{inflection} normally according to \isi{feminine} paradigms\is{paradigm}. There are still \textsc{M-forms}\is{M-form} (for `\isi{masculine}'), \isi{masculine} personal nouns with \isi{inflection} according to \isi{masculine} paradigms\is{paradigm}, but used with reference to women.\footnote{A non-negligible \isi{number} of words referring to women (which we could label \textsc{F-forms}) fall beyond the scope of the division presented here. These are inflectable\is{inflection} \isi{feminine} nouns that are morphologically independent from their \isi{masculine} counterparts, e.g. \emph{siostra} `sister' vs. \emph{brat} `brother'.} We should stress that there is actual competition among these word types, i.e. one lexical concept may be expressed by more than one type. (\ref{11ex9}--\ref{11ex10}) exemplify U-forms\is{U-form}.

\begin{exe}
	\ex \label{11ex9}
	\gll Z psycholog nie chciała gadać. \\
	with psychologist[\textsc{ins.sg}] \textsc{neg} want.\textsc{pst.f} talk.\textsc{inf} \\
	\trans ‘With the psychologist she wouldn’t talk.’ \hfill (NKJP)\footnote{Examples like (\ref{11ex9}) should theoretically be retrievable from corpora because of the charateristic mismatch between the preposition (taking an oblique \isi{case}) and the following \isi{U-form} which looks like a \isi{nominative}. Yet, such queries mostly yield noise, phrases with various forms wrongly tagged as nominatives\is{nominative}.}
	
	\ex \label{11ex10}
	\gll Zainteresowałem pomysł-em pani-ą notariusz Agnieszk-ę Sienkiewicz. \\
	interest.\textsc{pst.m.1sg} idea-\textsc{ins.sg} Mrs-\textsc{acc.sg} notary[\textsc{acc.sg}] Agnieszka-\textsc{acc.sg} Sienkiewicz[\textsc{acc.sg}] \\
	\trans `I got the notary, Agnieszka Sienkiewicz, interested in the idea.' \hfill (NKJP)
\end{exe}

{\sloppy The common nouns \emph{psycholog} in (\ref{11ex9}) and \emph{notariusz} in (\ref{11ex10}) as well as the surname in (\ref{11ex10}) have no marking even though their syntactic positions require oblique cases\is{case} that normally have non-zero suffixes. The three nouns are identical in form to the \isi{nominative} singular of their \isi{masculine} counterparts. With \isi{masculine} referents they would inflect\is{inflection}: \emph{z psychologi-em}, \emph{notariusz-a}, \emph{Sienkiewicz-a}, cf. the difference between (\ref{11ex4}) and (\ref{11ex5}). Applying the originally \isi{masculine} nouns for \isi{feminine} referents whilst keeping the unchanged basic form across the whole \isi{paradigm} is the formation of \textit{derived U-forms}\is{U-form}, a productive pattern of word-formation. The common nouns involved are usually ultimately borrowings, as in (\ref{11ex9}--\ref{11ex10}), but many are old ones and not necessarily perceived as foreign. Among surnames a great \isi{number} of U-forms\is{U-form} are native. Apart from derived U-forms\is{U-form}, there are also uninflectable\is{uninflectedness} loanwords referring to women, e.g. \emph{bizneswoman} in (\ref{11ex4}). These are not identical to their \isi{masculine} counterparts (or may not have them at all) and will be referred to as \textit{independent U-forms}\is{U-form}.\par}

To explain the \isi{uninflectedness} here, we need only generalise the analysis of (\ref{11ex4}). Words like \emph{psycholog}, \emph{notariusz}, etc. pose a problem for the Polish system if they are to refer to women. The word-final consonants in the basic form are characteristic of \isi{masculine} paradigms\is{paradigm}, but the words are supposed to be grammatically \isi{feminine} to match actual reference. As a result, the \isi{feminine} declension cannot be applied because of the morphonetic structure of the words, and the \isi{masculine} personal one is not used either because of semantics.\footnote{Section \ref{11sec2.1} mentions a \isi{paradigm} for feminines\is{feminine} ending in consonants. It might then seem that there is no conflict after all. Alas, this declension class only has nouns with (historically) soft	word-final consonants, which is not the \isi{case} for the majority of the words that appear as U-forms\is{U-form}. It also contains mostly inanimate nouns and few for animals, e.g. \emph{mysz} `mouse'. Any personal nouns following the pattern are long obsolete, e.g. \emph{jątrew} ‘husband's brother's wife'.}

Uninflectedness\is{uninflectedness} is only one possible resolution of the conflict of morphonetic shape and \isi{gender}. Apart from \isi{inflection}, Polish also has rich derivational morphology. If the morphonetic shape of a word is problematic for the given use, a derivational mechanism can rectify it. This is how D-forms\is{D-form} come into existence, examples of which are (\ref{11ex11}--\ref{11ex13}).

\begin{exe}
	\ex \label{11ex11}
	\gll Apetyczn-a trzydziestolat-k-a o wyglądzi-e prowincjonaln-ej nauczyciel-k-i. \\
	attractive-\textsc{nom.sg.f} thirty.year-\textsc{nmlz-nom.sg} \textsc{prep} appearance-\textsc{loc.sg} provincial-\textsc{gen.sg.f} teacher-\textsc{f-gen.sg} \\
	\trans ‘An attractive 30-year-old looking like a provincial teacher.’ \hfill (NKJP)
\end{exe}

The word \emph{nauczycielki} in (\ref{11ex11}) is clearly distinct from the U-forms\is{U-form} above and can be regarded as a regular element of the Polish system; the \isi{inflection} is present and it is in line with semantics. Note the feminative suffix \emph{-k(a)}, also present in (\ref{11ex2c}), which turns the \isi{masculine} noun \emph{nauczyciel} `teacher' into the \isi{feminine} \emph{nauczycielka}, which undergoes regular \isi{inflection}.

Most D-forms\is{D-form} are created by suffixation, but some are conversions (zero derivations).\footnote{Derived U-forms\is{U-form} are also conversions -- from a major \isi{masculine} personal \isi{paradigm} to the ``fully syncretic" class of uninflectable\is{uninflectedness} nouns.} A conspicuous (and controversial, cf. Section \ref{11sec4.4}) example is \emph{ministra} in (\ref{11ex12}), with exactly the same stem as the \isi{masculine} \emph{minister}. The word is formed by direct application of a major \isi{feminine} \isi{paradigm} (including \emph{-a} for the \isi{nominative} singular). Thus, even though the basic form \emph{minister-Ø} has no matching \isi{feminine} \isi{paradigm} (as was the \isi{case} with all the U-forms\is{U-form} above), the stem \emph{minist(e)r-} can be matched with no problem. This comes at the price of having a distinct basic form of the word, clearly a change noticeable for the speakers.

\begin{exe}
	\ex \label{11ex12}
	\gll Now-a ministr-a kultur-y Szwecj-i ma dred-y. \\
	new-\textsc{nom.sg.f} minister-\textsc{nom.sg} culture-\textsc{gen.sg} Sweden-\textsc{gen.sg} have.\textsc{3sg} dreadlock-\textsc{acc.pl} \\
	\footnotetext{\url{https://wyborcza.pl/7,75410,24399632,nowa-ministra-kultury-szwecji-ma-dredy-i-palila-trawke.html} Accessed 24.06.2024.\label{11ftn10}}
	\trans ‘The new Swedish minister for culture has dreadlocks.’\footnotemark
\end{exe}

The noun \emph{radna} `councillor, \textsc{f}' in (\ref{11ex13}) shares the stem with the \isi{masculine} \emph{radny}, but these are nouns of the adjectival \isi{paradigm}. Adjectives have \isi{gender} as an inflectional\is{inflection} category. Hence, whenever a reference to a woman needs to be made with concepts expressed by such nouns, the proper inflectional\is{inflection} suffixes are directly available in the \isi{paradigm}.

\begin{exe}
	\ex \label{11ex13}
	\gll Odpowiedź-Ø nie usatysfakcjonowała radn-ej Binder. \\
	answer-\textsc{nom.sg} \textsc{neg} satisfy.\textsc{pst.f} councillor-\textsc{gen.sg.f} Binder[\textsc{gen.sg}] \\
	\trans ‘The answer did not satisfy councillor Binder.’ \hfill (NKJP)
\end{exe}

Also the word \emph{trzydziestolatka} `30-year-old, \textsc{f}' in (\ref{11ex11}) has the same stem as the \isi{masculine} \emph{trzydziestolatek}. The two are parallel derivations from the phrase meaning `thirty years' rather than a direct instance of conversion. Still, the morphological closeness of both words is absolutely clear.

M-forms\is{M-form} are exemplified by the words \emph{nauczycielem} in (\ref{11ex14}), \emph{ministra} in (\ref{11ex15}) and \emph{autora} in (\ref{11ex1c}). Even though the reference is clearly \isi{feminine}, these are inflected\is{inflection} forms with suffixes of \isi{masculine} paradigms\is{paradigm}. The adjective in (\ref{11ex14}) has a \isi{masculine} suffix, which shows that M-forms\is{M-form} are grammatically \isi{masculine} nouns, just with female referents.

\begin{exe}
	\ex \label{11ex14}
	\gll Pan-i Grażyn-a nie jest surow-ym nauczyciel-em. \\
	Mrs-\textsc{nom.sg} Grażyna-\textsc{nom.sg} \textsc{neg} be.\textsc{3sg} strict-\textsc{ins.sg.m} teacher-\textsc{ins.sg} \\
	\trans ‘[Mrs] Grażyna [given name] is not a strict teacher.’ \hfill (NKJP)
	
	\ex \label{11ex15}
	\gll rocznic-a śmierc-i Izabel-i Jarug-i Nowack-iej, ministr-a polityk-i społeczn-ej \\
	anniversary-\textsc{nom.sg} death-\textsc{gen.sg} Izabela-\textsc{gen.sg} Jaruga-\textsc{gen.sg} Nowacka-\textsc{gen.sg.f} minister-\textsc{gen.sg} policy-\textsc{gen.sg} social-\textsc{gen.sg.f} \\
	\trans ‘the anniversary of the death of Izabela Jaruga-Nowacka, the Minister for Social Policy’\footnote{\url{https://www.gov.pl/web/rodzina/8-rocznica-smierci-izabeli-jarugi-nowackiej} \mbox{Accessed 24.06.2024.}}
\end{exe}

Let us add that (\ref{11ex14}) involves a noun used predicatively in a copular construction and this is one context in which M-forms\is{M-form} appear regularly and U-forms\is{U-form} are very rare. D-forms\is{D-form}, by contrast, are fine; an alternative way to express the meaning of (\ref{11ex14}) would be with the phrase \emph{surow-ą nauczyciel-k-ą}.

\begin{table}[ht]
	\centering
	\begin{tabularx}{\textwidth}{QQQQQ}%L{2.1cm}L{2.7cm}L{2cm}L{2cm}L{2.3cm}
		\lsptoprule
			& U-forms\is{U-form} & D-forms\is{D-form} & M-forms\is{M-form} & masculine\is{masculine} nouns \\
			\midrule
			reference & feminine\is{feminine} & feminine & feminine & masculine \\
			basic form (\textsc{nom sg}) & same as \textsc{m} noun (and M-form) & distinctly \textsc{f} & same as \textsc{m} noun (and U-form) & distinctly \textsc{m} (and same as U-form and M-form) \\
			inflectional\is{inflection} paradigm(s)\is{paradigm} & no marking & feminine & masculine & masculine \\
			morphological relation to \textsc{m} counterpart & identical to \textsc{m nom sg} (conversion) or no direct relation & suffixation or  (pseudo) conversion & full identity & n/a  \\
			examples in this section & \emph{psycholog} in (\ref{11ex9}), \emph{notariusz} in (\ref{11ex10}), \emph{Sienkiewicz} in (\ref{11ex10}) & \emph{panią} in (\ref{11ex10}), \emph{trzydziestolatka} in (\ref{11ex11}), \emph{nauczycielki} in (\ref{11ex11}), \emph{ministra} in (\ref{11ex12}), \emph{radna} in (\ref{11ex13}), \emph{Nowackiej} in (\ref{11ex15}) & \emph{nauczycielem} in (\ref{11ex14}), \emph{ministra} in (\ref{11ex15}) & n/a \\
		\lspbottomrule
	\end{tabularx}
	\caption{Summary of the properties of the discussed word types}
	\label{11tab2}
\end{table}

Table \ref{11tab2} summarises the properties of the three relevant word types. The distribution of the structures is partially determined by constructional, lexical and stylistic factors, discussed at length in Section \ref{11sec3.3}. The restriction on U-forms\is{U-form} being used predicatively has just been mentioned. As for the lexical factors, different concepts applied to women favour different structures, to the extent that only one of the three theoretically possible formations may be used. The role of register (and possibly ideology) can be noted with the words meaning `minister' above. (\ref{11ex12}) uses a \isi{D-form}, (\ref{11ex15}) an \isi{M-form}, a \isi{U-form} is also possible. (\ref{11ex12}) is part of a headline in a progressive daily \emph{Gazeta Wyborcza}, which at some point decided to adopt the innovative \isi{D-form} for female holders of ministerial offices, cf. Section \ref{11sec4.4}. (\ref{11ex15}) comes from a \isi{government} website and the \isi{M-form} has a dry bureaucratic feel to it. In fact, the \isi{U-form} \emph{minister}, as in (\ref{11ex24}), is the most common way to refer to female ministers.

Thus, there is some form of ``division of labour'' between the word types, but at the same time there is also competition between them. Even though we can clearly distinguish three types of noun usage, the speakers seem to perceive a binary opposition: D-forms\is{D-form} (showing \isi{feminine} reference directly) vs. the other two. The same transpires from linguistic analyses: even if the problem of \isi{uninflectedness} is acknowledged and discussed, it is subordinate to the main division: D-forms\is{D-form} are typically labelled \textit{\isi{feminine} forms} while U-forms\is{U-form} and M-forms\is{M-form} are \textit{\isi{masculine}}.\footnote{A notable exception is \citet[21, 172ff.]{krzyzanowski2013}.}

Let us also consider the status of the three word types as lexical units\is{lexeme}. It is obvious that D-forms\is{D-form} are separate lexemes\is{lexeme} from \isi{masculine} nouns. Likewise, it is clear that M-forms\is{M-form} are only special ways of using \isi{masculine} nouns. Independent U-forms\is{U-form}, like \emph{bizneswoman}, must also be lexemes\is{lexeme} of their own. The question is more complex for derived U-forms\is{U-form}. It is often assumed that they are twin lexemes\is{lexeme} of their \isi{masculine} counterparts, i.e. that there is an inflectable\is{inflection} \isi{masculine} \isi{lexeme} \emph{minister}\textsubscript{1} and a separate uninflectable\is{uninflectedness} \isi{feminine} one \emph{minister}\textsubscript{2}. The idea was apparently first suggested by \citet[68]{kucala1978} and is included in academic grammars (e.g. \citealt[368]{grzegorczykowapuzynina1998-problemy}; \citealt[422]{grzegorczykowa1998-rzeczownik}). Some more recent works refer to this position as a given (\citealt[126]{jadacka2006}; \citealt{rjp2012}; \citealt[626]{andrejewicz2019}), but it is
not universally accepted. \citet[227]{lazinski2006} argues strongly against it while \citet[229]{bobrowski2012} considers it only as a possible option. The issue will resurface in Sections \ref{11sec3.3.2} and \ref{11sec4.5}.

\begin{table}[ht]
	\centering
	\begin{tabularx}{\textwidth}{QQlQ}%L{2.5cm}L{2.5cm}L{2.2cm}L{2.8cm}
		\lsptoprule
			surname type 				      & masculine\is{masculine} \newline form & feminine\is{feminine} form & traditional feminine forms \\
											  & (further exs. in brackets) 			  & & (as of early 20\textsuperscript{th} cent. in standard language) \\
			\midrule
			suffix \emph{-sk(i)} and variants & Kowalski (Nowacki) 					  & Kowalska & Kowalska \\
			final \emph{-a} 				  & Duda (Jaruga) 						  & Duda & Dudzina/Dudowa (married) Dudzianka/Dudówna (unmarried) \\
			final consonant 				  & Nowak	(Sienkiewicz) 				  & Nowak (U-form\is{U-form}) & Nowakowa (married) Nowakówna (unmarried) \\
		\lspbottomrule
	\end{tabularx}
	\caption{Major inflectional\is{inflection} types of Polish surnames}
	\label{11tab3}
\end{table}

% \largerpage
Women's surnames present a similar level of variation to common nouns. As shown in Table \ref{11tab3}, surnames with the adjectival suffix \emph{-sk(i)} have regular D-forms\is{D-form}, while surnames with final consonants regularly have U-forms\is{U-form}. The most interesting group, not fitting any of the word types above, are those with the final \emph{-a}, which have identical inflectable\is{inflection} forms for men and women, with exactly the same \isi{inflection} in the singular, following the \isi{feminine} \isi{paradigm}. The grammatical \isi{gender} of the surname is decided by semantics across all the types. Traditionally, all types had distinct D-forms\is{D-form}, revealing not only \isi{gender}, but also a given woman's marital status.\footnote{There used to be greater variation of D-forms\is{D-form} than shown in Table \ref{11tab3}. Women could also be referred to by D-forms\is{D-form} based on their husbands or fathers' given names or positions. In the early modern times, similar derivative labels existed for young men, many of which developed into patronymic surnames.}

\subsection{U-forms\is{U-form} in use} \label{11sec3.2}

Let us explore the current usage of U-forms\is{U-form} in detail. (\ref{11ex16}--\ref{11ex17}) are quite typical examples.

\begin{exe}
	\ex \label{11ex16}
	\gll Głos-Ø Wybrzeż-a ma now-ą redaktor naczeln-ą. \\
	voice-\textsc{nom.sg} coast-\textsc{gen.sg} have.\textsc{3sg} new-\textsc{acc.sg.f} editor[\textsc{acc.sg}] main-\textsc{acc.sg.f} \\
	\trans ‘[The newspaper] Głos Wybrzeża has a new editor-in-chief.’ \hfill (NKJP)
	
	\ex \label{11ex17}
	\gll Lekarz powiedziała, że może bez problem-u jeść. \\
	doctor[\textsc{nom.sg}] say.\textsc{pst.f} that may.\textsc{3sg} without problem-\textsc{gen.sg} eat.\textsc{inf} \\
	\trans ‘The doctor said that [the baby] may eat [this kind of food] without a problem.’ \hfill (NKJP)
\end{exe}

The suffixes on noun modifiers as in (\ref{11ex16}) and on past verb forms in (\ref{11ex17}) are standard signals of the \isi{feminine} \isi{gender} and the U-forms\is{U-form} in the examples are clearly \isi{feminine}. However, such elements are not always present, cf. (\ref{11ex9}), where there is no overt marking of \isi{gender}, but the \isi{feminine} reference is clear too. This proves that among personal nouns the lack of inflection\is{uninflectedness} functions as a signal of \isi{feminine} \isi{gender} in its own right, as noted by \citet{obrebska-jablonska1949}.\footnote{Curiously	enough, even though she recognises \isi{uninflectedness} as a \isi{gender} signal, \citeauthor{obrebska-jablonska1949} herself describes the use of U-forms\is{U-form} as ``masculinisation'' (\citeyear[1]{obrebska-jablonska1949}), in line with the general tendency to block U-forms\is{U-form} and M-forms\is{M-form} together against D-forms\is{D-form}, cf. Section \ref{11sec3.1}.}

If a \isi{U-form} appears in the \isi{subject} position, it is identical to the corresponding \isi{masculine} noun with a zero suffix. A past tense verb may show \isi{gender} as in (\ref{11ex17}), but not all sentences are past. Those like (\ref{11ex18}) leave the \isi{gender} of the referent underdetermined. This follows directly from the rules that establish the class of U-forms\is{U-form}, but is untypical of Polish, esp. with specific referents \citep[263]{lazinski2006}, cf. fn. \ref{11ftn17} and the discussion of (\ref{11ex28}) and (\ref{11ex32}).

\begin{exe}
	\ex \label{11ex18}
	\gll Jakby nie rozumiał, co doktor mówi. \\
	as\_if \textsc{neg} understand.\textsc{pst.m} what.\textsc{acc} doctor[\textsc{nom.sg}] say.\textsc{3sg} \\
	\trans ‘As if he didn’t understand what the doctor was saying.’ \hfill (NKJP)\footnote{It is impossible to retrieve easily examples of U-forms\is{U-form} like (\ref{11ex18}) from corpora. U-forms\is{U-form} are tagged the same as \isi{masculine} nouns and true examples of the latter dominate the results absolutely.}
\end{exe}

Misunderstanding sentences like (\ref{11ex18}) is probably relatively seldom a risk in real-life communication as the \isi{gender} will have been established earlier in the discourse.\footnote{In the context of (\ref{11ex18}), we learn two sentences earlier that the doctor is a woman from her \isi{D-form} surname.\label{11ftn17}} Without any such cues the default interpretation would likely be that the sentence is about a male doctor.\footnote{We can then see a point in regarding U-forms\is{U-form} as ``\isi{masculine}''.} Thus, plain U-forms\is{U-form}, even though attested and used, may feel somewhat inadequate.

In fact, U-forms\is{U-form} are very often combined in \isi{apposition} with other \isi{feminine} nouns, typically inflectable\is{inflection} ones. In (\ref{11ex19}--\ref{11ex20}) U-forms\is{U-form} are preceded by more general inflectable\is{inflection} nouns, which directly show the syntactic relation of the phrases to the surrounding context. The \isi{uninflectedness} becomes much less problematic in \isi{apposition}, the \isi{U-form} functions like a label quoted in an unchanged basic form, as in (\ref{11ex6}), (\ref{11ex7a}) and (\ref{11ex8}).

\begin{exe}
	\ex \label{11ex19}
	\gll T-a spraw-a jest ważn-a dla pan-i prezes. \\
	this-\textsc{nom.sg.f} matter-\textsc{nom.sg} be.\textsc{3sg} important-\textsc{nom.sg.f} for Mrs-\textsc{gen.sg} president[\textsc{gen.sg}] \\
	\trans ‘This matter is important for [you,] madame president [of the company].’
	
	\hfill (NKJP)
	
	\ex \label{11ex20}
	\gll Węgr-y też mają kobiet-ę minister spraw-Ø zagraniczn-ych. \\
	Hungary-\textsc{nom.pl} also have.\textsc{3pl} woman-\textsc{acc.sg} minister[\textsc{acc.sg}] affair-\textsc{gen.pl} foreign-\textsc{gen.pl} \\
	\trans ‘Hungary also has a female foreign minister.’ \hfill (NKJP)
\end{exe}

The use of the honorific \emph{pani} in \isi{apposition} with U-forms\is{U-form}, as in (\ref{11ex19}) or (\ref{11ex10}) above, is a regular pattern in the language. Other words can be used similarly, but this rare usage always meets specific needs of a given context, e.g. in (\ref{11ex20}) the noun \emph{kobietę} explicitly labels the \isi{gender} of the referent. The honorific combines regularly not only with common nouns, but also with surnames of women, both inflectable\is{inflection} and uninflected\is{uninflectedness}. Table~\ref{11tab4} shows the importance of the appositional\is{apposition} pattern in terms of frequency. With the phrase \emph{dla pani prezes} from (\ref{11ex19}) as a starting point, the NKJP was queried for the preposition \emph{dla} and the noun form \emph{prezes} with one intervening word possible. The results show a strong dominance of the \isi{apposition} with the honorific for the \isi{U-form}. No such effect is visible in the data for corresponding \isi{masculine} phrases. This suggests that \emph{pani} has developed into a kind of auxiliary noun with U-forms\is{U-form}, a carrier of overt inflectional\is{inflection} marking and analytic marker of the \isi{feminine} \isi{gender}, cf. the discussion of (\ref{11ex32}). Note also that in the query results there is no single example parallel to (\ref{11ex20}), i.e. with a common noun other than \emph{pani} used in \isi{apposition} with the noun \emph{prezes}.

\begin{table}[ht]
	\centering
	\begin{tabularx}{0.9\textwidth}{L{7cm}rr}
		\lsptoprule
			phrase structure & no. of relevant & \% \\
			& examples & \\
			\midrule
			\emph{dla pani prezes} (honorific + U-form\is{U-form}) & 43 & 80\% \\
			\emph{dla} \textsc{adj.f/det.f} \emph{prezes} (\textsc{f} modifier + U-form) & 4 & 7\% \\
			\emph{dla prezes} (no intervening words) & 7 & 13\% \\
			\midrule
			total \textsc{f} & 54 & \\
			\midrule
			& & \\
			\emph{dla pana prezesa} (honorific + \textsc{m} noun) & 74 & 6\% \\
			\emph{dla} \textsc{adj.m/det.m} \emph{prezesa} (modifier+ \textsc{m} noun) & 137 & 11\% \\
			\emph{dla prezesa} (no intervening words) & 1,034 & 83\% \\
			\midrule
			total \textsc{m} & 1,245 & \\
		\lspbottomrule
	\end{tabularx}
	\caption{Frequencies of various types of phrase structure in phrases \emph{dla (...) prezes(a)} ‘for the president (of a company)’ in the NKJP.}
	\label{11tab4}
\end{table}

As for women's surnames, they commonly appear in \isi{apposition}, with the noun \emph{pani} and with given names, most of which inflect\is{inflection}.\footnote{The vast majority of female given names in Polish end in \emph{-a} in the basic form and inflect\is{inflection} according to the major \isi{feminine} paradigms\is{paradigm}. The few rather rare exceptions are U-forms\is{U-form}.\label{11ftn18}} In (\ref{11ex10}), such typical structures are joined into a single complex \isi{apposition}: honorific -- profession -- given name -- surname. Two nouns are inflected\is{inflection}, two are U-forms\is{U-form}.

Another natural combination is that of a given woman's function and her surname, as in (\ref{11ex13}). In fact, of the 7 examples with no intervening word before the \isi{U-form} in Table \ref{11tab4}, as many as 4 involve a surname in \isi{apposition}, coming after the word \emph{prezes}. In such phrases both nouns may well be U-forms\is{U-form}, as in (\ref{11ex21}).

\begin{exe}
	\ex \label{11ex21}
	\gll polityczn-a burz-a po twitterow-ym wpisi-e kurator Nowak \\
	political-\textsc{nom.sg.f} storm-\textsc{nom.sg} after twitter.\textsc{adj-loc.sg.m} post-\textsc{loc.sg} superintendent[\textsc{gen.sg}] Nowak[\textsc{gen.sg}] \\
	\trans ‘a political storm after a tweet by the school superintendent Nowak’\footnote{\url{https://wiadomosci.dziennik.pl/polityka/artykuly/8578116,barbara-nowak-lex-czarnek-barbara-nowacka.html} Accessed 26.06.2024.}
\end{exe}

As in (\ref{11ex9}), there are no direct signals of \isi{feminine} \isi{gender} in (\ref{11ex21}), but the \isi{feminine} reference is perfectly clear because of the lack of expected inflection\is{uninflectedness}. A parallel phrase about a male superintendent by the same name would obviously need to inflect\is{inflection}: \emph{kurator-a Nowak-a}, cf. the discussion of (\ref{11ex4}--\ref{11ex5}) and (\ref{11ex9}--\ref{11ex10}).

The opposition between the inflectable\is{inflection} male version of a surname and the uninflected\is{uninflectedness} female one extends to many foreign surnames. For example, the surname \emph{Clinton} inflects\is{inflection} when Bill Clinton is referred to, but does not when Hillary Clinton is, but cf. Section \ref{11sec4.2}. Here the given names behave identically to the surname: \emph{Bill} is regularly inflected\is{inflection} as a \isi{masculine} noun, while \emph{Hillary}, not fitting any \isi{feminine} \isi{paradigm}, is a \isi{U-form}.

(\ref{11ex22}), the title of a paper devoted to U-forms\is{U-form}, exemplifies how surnames function as plain U-forms\is{U-form}. As we can see, this is parallel to what we observed with common nouns.

\begin{exe}
	\ex \label{11ex22}
	\gll Baran mówi o Kowal. \\
	Baran\textsc{[nom.sg]} talk.\textsc{3sg} about Kowal\textsc{[loc.sg]} \\
	\trans ‘Baran talks about [Ms.] Kowal.’ \hfill \citep{pawlowski1951}
\end{exe}

The surname \emph{Kowal} in (\ref{11ex22}) is in a non-subject position and the lack of overt inflection\is{uninflectedness} identifies it as a \isi{U-form} and a woman's name, as in (\ref{11ex9}) and (\ref{11ex21}). However, for the surname \emph{Baran}, which appears as the \isi{subject}, it cannot be decided from the sentence alone if the person in question is a man or a woman, as in (\ref{11ex18}).

To sum up, U-forms\is{U-form} are an established and productive pattern. Among them there are common nouns for professions, professional functions and titles as well as surnames (including numerous native ones). These words do not have any overt marking of the \isi{feminine} \isi{gender}, but trigger \isi{feminine} \isi{agreement}. However, the lack of inflection\is{uninflectedness} on nouns is a conspicuous absence in Polish\footnote{\citet[35]{krzyzanowski2013} vividly describes uninflected\is{uninflectedness} nouns, including U-forms\is{U-form}, as ``refraining from inflection\is{uninflectedness}''.} and so it has developed into a signal of the \isi{feminine} \isi{gender} by itself. This signalling fails, however, if the given noun is in the \isi{nominative} singular, which for derived U-forms\is{U-form} is identical to the \isi{masculine} counterpart. There is no grammatical difference between common nouns and personal names. All the effects related to \isi{uninflectedness} can be observed in both groups.

The \isi{feminine} \isi{agreement} triggered by U-forms\is{U-form} is the reason \citet[61, 68]{kucala1978} concluded they are \isi{feminine} nouns, thus excluding any mismatch of \isi{gender} between the modifiers or past verb forms and the nouns of an apparently \isi{masculine} shape.

Despite their established status in the system, U-forms\is{U-form} are still a very odd class in Polish. This fuels the use of a rival pattern, D-forms\is{D-form}. Taking into account the most general features of the language, one might expect a very strong systemic pressure for establishing D-forms\is{D-form} across the whole personal noun lexicon. As can be inferred so far, this has not been the \isi{case}, but see the discussion of recent developments in Section \ref{11sec4.4}.

Another effect related to the ``odd'' character of U-forms\is{U-form} is the ease with which they enter into appositional\is{apposition} structures. In this way, they can ``hide" behind inflectable\is{inflection} nouns, which show syntactic relations through \isi{case} marking and manifest \isi{gender} explicitly. The main inflectable\is{inflection} nouns entering into these structures are the honorific \emph{pani} and given names. Apposition\is{apposition} is thus a major strategy of integrating U-forms\is{U-form} in the grammatical system heavily reliant on \isi{inflection}.

\is{U-form|(}
\subsection{Constraints on the use of U-forms} \label{11sec3.3}

\subsubsection{Overview}

The use of U-forms is limited in several ways in present-day Polish. Some restrictions are grammatical. One is the virtual absence of U-forms from the copular construction where M-forms\is{M-form} are used rather than U-forms, cf. the discussion of (\ref{11ex14}), which need not be expanded. U-forms are also less free in the plural, see the next section. Apart from that there are lexical and stylistic factors. Groups of lexemes\is{lexeme} show preferences for particular word types -- this is the ``division of labour'' between U-forms, D-forms\is{D-form} and M-forms\is{M-form}. However, there is also substantial (and increasing) overlap and competition between the types. When forms of two or even three types co-exist for the same lexical concept, the difference between them is related to factors like register or ideology.

\subsubsection{Plural \isi{number}} \label{11sec3.3.2}

\largerpage
The examples so far have all shown U-forms in the singular. Number\is{number} is the most problematic category for U-forms. As has been shown in Section \ref{11sec3.2}, \isi{case} can usually be inferred from the context as its value must match the requirements of the word governing a given U-form syntactically, even though structurally ambiguous phrases appear in actual use cf. \citet[340--341]{szpyra-kozlowska2021}. Gender\is{gender} is understood to be \isi{feminine} due to the lack of inflection\is{uninflectedness}. A U-form by itself has no way of differentiation left for \isi{number}, which is by default understood to be singular.\footnote{An alternative would be for the \isi{number} to be underspecified in such contexts. This is not the \isi{case}, however, e.g. (\ref{11ex9}) refers to one female psychologist, a plural reading is highly unlikely if not downright impossible.} A natural solution is to mark plurality around the noun: on nouns in \isi{apposition}, modifiers or verbs. In (\ref{11ex23}) the U-form surname is combined with the plural of the honorific \emph{pani}. Yet again \isi{apposition} makes the functioning of an uninflected\is{uninflectedness} noun in Polish syntax possible.

\begin{exe}
	\ex \label{11ex23}
	\gll Postanowiono pomóc pani-om Anflik w inn-y sposób-Ø. \\
	decide.\textsc{pst.impers} help.\textsc{inf} Mrs-\textsc{dat.pl} Anflik[\textsc{dat.pl}] in other-\textsc{acc.sg.m.inan} way-\textsc{acc.sg} \\
	\trans ‘It was decided that the Anflik Ladies will be helped in a different way.’
	
	\hfill (NKJP)
\end{exe}

The plural honorific is used regularly, but other potential carriers of plurality are not unproblematic. \citet[227]{lazinski2006} argues that derived U-forms simply cannot be used as plural nouns in their own right, giving an example with the ``evidently faulty'' phrase \emph{*między dwiema minister} `between two [female] ministers'. A phrase with an \isi{apposition}, parallel to (\ref{11ex23}), \emph{między paniami minister} would be acceptable, however. An independent U-form in plural is also fine: \emph{między dwiema lady} `between two ladies'.

Łaziński's conclusion is that derived U-forms are not separate uninflectable\is{uninflectedness} lexemes\is{lexeme}, but simply special cases\is{case} of usage of \isi{masculine} nouns (otherwise normally inflectable\is{inflection}) for female referents. The lack of marking would then prove to be simply contextual \isi{uninflectedness} -- in contrast to the view taken by \citet{kucala1978}, cf. Section \ref{11sec3.1}. It follows that there are two separate ways in which \isi{masculine} personal lexemes\is{lexeme} can be used to refer to women, U-forms and M-forms\is{M-form}. The analysis runs into problems when we consider that U-forms trigger \isi{feminine} \isi{agreement}. \citet[231]{lazinski2006} explains the resultant \isi{gender} mismatch (non-existent in Kucała's approach, cf. Section \ref{11sec3.2}) as ``\isi{agreement} \emph{ad sensum}'', i.e. a grammatically \isi{masculine} noun combines with \isi{feminine} modifiers or past verb forms on the basis of actual \isi{feminine} reference. This special kind of \isi{agreement} would then be obligatory for U-forms, but never happens with M-forms\is{M-form}, cf. (\ref{11ex14}). The firmness of Łaziński's conclusions is shaken by the fact that sentences with the supposedly unacceptable structures are in fact attested, e.g. (\ref{11ex24}).

\begin{exe}
	\ex \label{11ex24}
	\gll Słow-a t-e kierowały obi-e minister wyłącznie do osób-Ø, \mbox{któr-ych (...)} \\
	word-\textsc{acc.pl} this-\textsc{acc.pl.nvir} direct.\textsc{pst.pl.nvir} both-\textsc{nom.pl.f} minister[\textsc{nom.pl}] only to person-\textsc{gen.pl} who-\textsc{gen.pl} \\
	\trans ‘These words were spoken by both [female] ministers only to people \mbox{who (...)}’ \hfill (NKJP)
\end{exe}

In (\ref{11ex24}) plurality is marked on the modifier and the verb form. Ironically, the example includes the very same noun as Łaziński's ``faulty'' phrase. Yet, sentences like (\ref{11ex24}) are rather infrequent. To be fair, Łaziński gives a counterexample himself, quoting what he had heard said by a receptionist at the doctor's, given as (\ref{11ex25}).

\begin{exe}
	\ex \label{11ex25}
	\gll Doktor już przyszły. \\
	doctor[\textsc{nom.pl}] already come.\textsc{pst.pl.nvir} \\
	\trans ‘The doctors [all female] have come.’ \hfill (\citealt[227]{lazinski2006})
\end{exe}

Plurality in (\ref{11ex25}) is shown on the verb only. Łaziński comments that this kind of usage ``does not seem to be very common". This observation seems accurate, but such structures still do occur in real language. (\ref{11ex26}) is a parallel example with an independent U-form. Such sentences are easier to find than ones like (\ref{11ex25}) even though independent U-forms are generally less common than derived U-forms.

\begin{exe}
	\ex \label{11ex26}
	\gll Dawniej miss były zdan-e tylko na siebie. \\
	early.\textsc{comp.adv} beauty\_queen[\textsc{nom.pl}] be.\textsc{pst.pl.nvir} rely.\textsc{ptcp-nom.pl.nvir} only on oneself.\textsc{acc} \\
	\trans ‘In the earlier times, beauty queens were left on their own.’\footnote{ \url{https://dziennikzachodni.pl/ewa-wojciechowska-dawniej-miss-byly-zdane-tylko-na-siebie-to-byl-blad/ar/486300} Accessed 26.06.2024.}
\end{exe}

In the light of (\ref{11ex24}--\ref{11ex26}) it is impossible to accept Łaziński's categorical statements at face value, but the observations are not irrelevant; the counterexamples are infrequent. Table \ref{11tab5} sums up the use of U-forms in plural. The possible carriers of plurality form a hierarchy. Let us add that we are not aware of attestations of U-forms with clearly plural meaning and no formal marking of plurality. If such usage exists at all, it is rare, quite certainly substandard and strongly context-dependent. In contrast to U-forms, the other patterns: D-forms\is{D-form} and M-forms\is{M-form} as well as \isi{masculine} personal nouns are used in the plural without restrictions.

\begin{table}[ht]
	\centering
	\begin{tabularx}{0.85\textwidth}{L{4cm}L{3.5cm}L{1.5cm}}
		\lsptoprule
			locus of plural marking & occurrence & examples \\
			\midrule
			\emph{pani} in apposition\is{apposition} & regular & (\ref{11ex23}) \\
			noun modifiers & attested but not common & (\ref{11ex24}) \\
			finite verb & rare, mostly with independent U-forms & (\ref{11ex25}--\ref{11ex26}) \\
		\lspbottomrule
	\end{tabularx}
	\caption{Occurrence of U-forms in plural contexts}
	\label{11tab5}
\end{table}

In Table \ref{11tab5} we present the general picture of the usage of U-forms in the plural. Since it is extremely difficult to substantiate the frequency labels for the whole class, we resort to corpus data for specific example words, shown in Table \ref{11tab6}. We selected three typical U-forms: \emph{minister}, \emph{doktor} `medical doctor, Ph.D.' and \emph{profesor} `professor'. The data have been queried separately for different loci of plural marking. These were: the non-nominative plural forms of the honorific, non-nominative \isi{feminine} forms of basic plural modifying words: \emph{dwie} `two', \emph{obie} `both', \emph{wsystkie} `all' and lastly top-frequency plural verb forms: \emph{są} `they are', \emph{będą} `they will (be)' and \emph{mają} `they have'. The first two are placed before the target U-forms in the query, the verb forms after them. The use of non-nominative forms in the first two excludes plural \isi{agreement} on the verb. For comparison, Table~\ref{11tab6} also includes figures for parallel plural phrases with a single \isi{D-form} \emph{nauczycielka} `teacher, \textsc{f}'.

\largerpage
\begin{table}[H]
	\centering
	\begin{tabularx}{\textwidth}{L{5cm}R{2.2cm}R{2.7cm}R{1cm}}
		\lsptoprule
		phrase structure & no. of query results & exs. with single locus of plurality & \% \\
		\midrule
		\emph{pani} (\textsc{pl, nnom}) + \emph{minister}/\emph{doktor}/\emph{profesor} & 23 & 19 & 95\% \\
		\emph{dwie}/\emph{obie}/\emph{wszystkie} (\textsc{pl, nnom}) + \emph{minister}/\emph{doktor}/\emph{profesor} & 1 & 1 & 5\% \\
		\emph{minister}/\emph{doktor}/\emph{profesor} + \emph{są}/\emph{będą}/\emph{mają} & 3 & 0 & 0\% \\
		\midrule
		total U-forms & & 20 & \\
		\midrule
		\emph{pani} (\textsc{pl, nnom}) + \emph{nauczycielki} (\textsc{pl, nnom}) & 83 & 81 & 47\% \\
		\emph{dwie}/\emph{obie}/\emph{wszystkie} (\textsc{pl, nnom}) + \emph{nauczycielki} (\textsc{pl, nnom}) & 41 & 41 & 24\% \\
		\emph{nauczycielki} + \emph{są}/\emph{będą}/\emph{mają} & 56 & 49 & 29\% \\
		\midrule
		total D-forms & & 171 & \\
		\lspbottomrule
	\end{tabularx}
	\caption{Frequencies of various types of plural phrase structure for selected U-forms and D-forms\is{D-form} in the NKJP}
	\label{11tab6}
\end{table}

% \needspace{2\baselineskip}
Since we are interested in comparing different loci of plural marking, examples with doubled signals of plurality (e.g. both on the modifier and the honorific) had to be removed from the final count. In the \isi{case} of U-forms combining with the selected plural verb forms this meant excluding all examples. We have only a single example, cited as (\ref{11ex24}), of a plural modifier with a U-form even though three different words and three possible modifiers were used in the query. Table~\ref{11tab6} shows both the initial figures and the examples with single locus of plurality after the recount. The data confirm the hierarchy of carriers of plural marking suggested in Table~\ref{11tab5}. The honorific is significantly more common relative to the other types in U-forms than in D-forms\is{D-form}. The plural honorific dominates in the results with the U-forms. The distribution of the structures with the \isi{D-form} is much more even.

\subsubsection{Lexical range} \label{11sec3.3.3}

Various areas of the personal noun lexicon of Polish show preferences for particular word types when \isi{feminine} reference is made. In a way, each of the three types is restricted in its lexical range; default or at least common for some words, but rare or even non-existent for some others. We will briefly outline what kind of nouns prefer each of the morphological types, starting with words dispreferring U-forms.

Very many concepts are regularly expressed by D-forms\is{D-form} when referring to women. (\ref{11ex27}) shows pairs of corresponding words for men and women, the \isi{feminine} ones show relevant morphological division.

\begin{exe}
	\ex \label{11ex27}
	\begin{xlist}
		\ex \label{11ex27a} Włoch {\sim} Włosz-k-a
		\trans ‘Italian person’
		
		\ex \label{11ex27b} kuzyn {\sim} kuzyn-k-a
		\trans ‘cousin’
		
		\ex \label{11ex27c} sąsiad {\sim} sąsiad-k-a
		\trans ‘neighbour’
		
		\ex \label{11ex27d} artysta {\sim} artyst-k-a
		\trans ‘artist’
		
		\ex \label{11ex27e} poseł {\sim} posł-ank-a
		\trans ‘MP’
		
		\ex \label{11ex27f} sprzedawca {\sim} sprzedawcz-yn-i
		\trans ‘shop assistant’
		
		\ex \label{11ex27g} szef {\sim} szef-ow-a
		\trans ‘boss’
		
		\ex\label{11ex27h} hrabia {\sim} hrab-in-a
		\trans ‘count(ess)’
		
		\ex\label{11ex27i} diabeł {\sim} diabl-ic-a
		\trans ‘devil, evil person’
		
		\ex \label{11ex27j} staruch {\sim} staruch-a
		\trans ‘old person (derog.)’
	\end{xlist}
\end{exe}


Word pairs such as those in (\ref{11ex27}) are most consistent with general rules of Polish and constitute the ``basic'' \citep[49]{kucala1978} section among personal nouns. Feminine\is{feminine} reference is made with D-forms\is{D-form} and U-forms are rare or non-existent. Most of these concepts allow M-forms\is{M-form} in copular constructions and \isi{masculine} nouns in generic use (covering both men and women).\footnote{(\ref{11ex2a}--\ref{11ex2b}) are examples of \isi{masculine} nouns used generically. An alternative to generic masculines\is{masculine} is splitting, as in (\ref{11ex2c}).}

As for semantics, there are names for nationalities in this group, cf. (\ref{11ex27a}), as well as family relations, cf. (\ref{11ex27b}). Apart from that, it is difficult to generalise, but we must note that some words for occupations and functions also appear here, cf. (\ref{11ex27d}--\ref{11ex27f}), as well as those for noble titles, cf. (\ref{11ex27h}).

The standard assumption in Polish linguistics is that D-forms\is{D-form} are derived from their \isi{masculine} counterparts (e.g. \citealt[422]{grzegorczykowa1998-rzeczownik}; \citealt[49]{kucala1978}). While this is right in many cases\is{case}, \citet[134--135]{szpyra-kozlowska2021} points out that it need not be a given. For example, nationality labels for both men and women, as in (\ref{11ex27a}), are most logically derived from country names, cf.~also the discussion of \emph{trzydziestolatka} in (\ref{11ex11}).

(\ref{11ex27}) shows the rich array of feminative suffixes in Polish: the most common \emph{-k(a)} in (\ref{11ex27a}--\ref{11ex27d}), as well as (\ref{11ex2c}) and (\ref{11ex11}) above, \emph{-ank(a)} in (\ref{11ex27e}), \emph{-yn(i)/}\emph{-in(i)} in (\ref{11ex27f}), \emph{-ow(a)} in (\ref{11ex27g}), \emph{-in(a)} in (\ref{11ex27h}), \emph{-ic(a)} in (\ref{11ex27i}). Finally, some D-forms\is{D-form} are conversions, cf. (\ref{11ex27j}) and (\ref{11ex12}).

The suffix \emph{-k(a)} is the most common and can be considered default. Its disadvantage in the eyes of the speakers is that it is identical to the diminutive one, and the one forming names of objects from personal nouns, e.g \emph{marynar-k-a }`suit jacket, \textsc{f}' from \emph{marynarz} `(navy) sailor, \textsc{m}', as well as the nominalising one, as in \emph{trzydziestolatka} in (\ref{11ex11}).\footnote{The identity of the feminative -\emph{k(a)} and suffixes with other meanings may lead to emergence of excessive homophones, which is a factor that supposedly blocks the creation of D-forms\is{D-form} in some cases\is{case}. \citet[191--192, 262--265]{szpyra-kozlowska2021}, however, points out that homophony of suffixes is common in Polish derivation.\label{11ftn25}} The variety of suffixes itself may also be problematic when new D-forms\is{D-form} are created as the speakers might not be sure which suffix to apply (cf. \citealt{szpyra-kozlowska2019}, see also Sections \ref{11sec4.3} and \ref{11sec4.4}).

As for \isi{inflection}, most D-forms\is{D-form} follow major \isi{feminine} paradigms\is{paradigm}, those in \emph{\nobreakdash-yn(i)}/\emph{-in(i)} belong to a minor class, while the suffix \emph{-ow(a)} builds words of the adjectival declension. Note that even if the \isi{masculine} noun has the typically \isi{feminine} \emph{-a} in the basic form, D-forms\is{D-form} are usually distinct, cf. (\ref{11ex27d}), (\ref{11ex27f}) and (\ref{11ex27h}).

At the other pole, there are lexemes\is{lexeme} that can refer to both men and women, but have forms inflected\is{inflection} according to \isi{masculine} paradigms\is{paradigm}, i.e. they are regularly M-forms\is{M-form} when the reference is \isi{feminine}. D-forms\is{D-form} and especially U-forms are rare for such words, but not completely non-existent, cf. (\ref{11ex30}) and Section \ref{11sec4.1}. Examples include \emph{szpieg} `spy', \emph{wróg} `enemy', \emph{przechodzień} `passer-by', \emph{tchórz} `coward' or \emph{gość} `guest' as in (\ref{11ex28}).

\begin{exe}
	\ex \label{11ex28}
	\gll Musi przywitać nasz-ego gości-a osobiście. \\
	must.\textsc{3sg} welcome.\textsc{inf} our-\textsc{acc.sg.m} guest-\textsc{acc.sg} personally \\
	\trans ‘He must welcome our guest in person.’ \hfill (NKJP)
\end{exe}

The noun \emph{gościa} in (\ref{11ex28}) takes a suffix from a \isi{masculine} \isi{paradigm} and the \isi{masculine} \isi{gender} is evident on the modifier. Without context the sentence may well refer to a male guest. In fact, this would be the default interpretation. A female name appears several paragraphs earlier in the text (\ref{11ex28}) is taken from, thus making the actual reference clearly \isi{feminine}.

\citet[267]{lazinski2006} claims that the nouns of this group are rarely used for specific reference, which is normally associated with clear marking of the referent's \isi{gender} in Polish. The mismatch between \isi{feminine} reference and ostensibly \isi{masculine} form of the nouns is mitigated by the backgrounding of \isi{gender} in their meaning and non-specific reference. For example, \emph{tchórz} focuses on the feature of character, while \emph{przechodzień} emphasises a momentary role in traffic. This line of reasoning fits some of the words less well, e.g. \emph{gość}, which may often have specific reference, as it does in (\ref{11ex28}). Similarly, the noun \emph{szpieg} could logically refer to specific individuals as well, the synonymous \emph{agent} is apparently more likely in such contexts \citep[268]{lazinski2006}. The latter has a regular \isi{D-form} \emph{agentka} for women as referents.

Apart from the fully explicit marking of femininity by D-forms\is{D-form} and its occlusion with M-forms\is{M-form}, there is a group of concepts for which U-forms are used regularly. For some of these words, this is the dominant pattern. D-forms\is{D-form} may exist as alternatives, but their prevalence and acceptability varies greatly depending on the word. The words appearing as derived U-forms so far in the text are mostly labels for professions and professional titles: \emph{notariusz} `notary', \emph{psycholog} `psychologist', \emph{redaktor} `editor', \emph{lekarz} `medical doctor', \emph{doktor} `doctor', \emph{prezes} `president (e.g. of a company)', \emph{kurator} `educational superintendent', \emph{profesor} `professor', \emph{minister} `minister'. What is common for these words is some level of prestige.\footnote{In the 20\textsuperscript{th} century Polish linguists suggested various factors other than prestige as blocking the creation of D-forms\is{D-form}, and promoting U-forms/M-forms\is{M-form} in the lexicon, cf. \citet[108--110]{buttlerkurkowskasatkiewicz1973}. \citet[190--241]{szpyra-kozlowska2021} offers a convincing criticism of these factors, calling them ``myths''. One of them is the identity of the feminative \emph{-k(a)} and suffixes with other meanings, cf. fn. \ref{11ftn25}.} Indeed, the more prestigious a concept is, the more likely it is to be expressed by a U-form. The model for such usage across various contexts is direct address, which takes the form of the \isi{apposition} comprising the honorific \emph{pani} and the given U-form, cf. (\ref{11ex19}). \citet[275--276]{lazinski2006} claims the use of a word as a regular title of address blocks the possibility of creating a \isi{D-form}. The blockage is not complete, and has been loosening recently, cf. Section \ref{11sec4.4}. However, the appositional\is{apposition} structure for addressing women remains a stronghold of U-forms. (\ref{11ex29}) is an example of a creative use of the pattern. The noun \emph{ekspert} is not normally used for direct address and has an established \isi{D-form} \emph{ekspertka}. However, in the given situation, the U-form was chosen as a way of semi-direct address.

\begin{exe}
	\ex \label{11ex29}
	\gll T-o dla pan-i ekspert. \\
	this-\textsc{nom.sg} for Mrs-\textsc{gen.sg} expert[\textsc{gen.sg}] \\
	\trans ‘This is for [you,] expert.’ \hfill (NKJP)
\end{exe}

{\sloppy (\ref{11ex30}--\ref{11ex31}) show words with relatively little prestige used as U-forms without the honorific. The examples are recent press headings which \citet{szpyra-kozlowska2021} presents as instances of blatant overuse of the uninflected\is{uninflectedness} pattern.\par}

\begin{exe}
	\ex \label{11ex30}
	\gll polsk-a szpieg uhonorowan-a w Londyni-e \\
	Polish-\textsc{nom.sg.f} spy[\textsc{nom.sg}] honour.\textsc{ptcp-nom.sg.f} in London-\textsc{loc.sg} \\
	\trans ‘A Polish [woman] spy [from WW2] honoured in London’
	
	\hfill \citep[339]{szpyra-kozlowska2021}
	
	\ex \label{11ex31}
	\gll Piękn-a sołtys wyznaje. \\
	beautiful-\textsc{nom.sg.f} village\_leader[\textsc{nom.sg}] confess.\textsc{3sg} \\
	\trans ‘The beautiful village leader confesses.’ \hfill (339)
\end{exe}

(\ref{11ex30}) features the use of \emph{szpieg} as a U-form, although the concept is usually expressed as an \isi{M-form}. For \emph{sołtys} in (\ref{11ex31}) a more or less established \isi{D-form} exists. The examples are certainly not typical for this word type and may seem questionable for some speakers. On the other hand, such examples of U-forms prove the pattern is productive.

In sum, the nouns preferring D-forms\is{D-form} are very much regular: their form and \isi{inflection} are in line with semantics. The ones preferring M-forms\is{M-form} also inflect\is{inflection}, but ignoring the actual \isi{feminine} reference. U-forms are exceptional because of their \isi{uninflectedness}. As can be seen, they are not distributed freely in the lexicon, but concentrate in the domain of words for professions and professional titles. Despite these divisions in the lexicon, there is still a considerable field for competition between the three types.

\subsubsection{Register} \label{11sec3.3.4}

When they are in competition, U-forms are somewhat formal while the competing D-forms\is{D-form} are informal and often colloquial. This matches the factor of prestige, cf. Section \ref{11sec3.3.3}. Let us consider the word \emph{dyrektor} `school principal' as an example. Since many principals in Polish schools are women, it is naturally used as a U-form, typically with the honorific: \emph{pani dyrektor}. This form is expected and used in contexts like teacher-student or teacher-parent communication, and esp.~when addressing the principals. However, in less formal contexts, e.g. between students or between fellow teachers, the \isi{D-form} \emph{dyrektorka} is used as well. This pattern is found with many U-forms. It is generally easier for the U-forms to appear in less formal contexts, cf. (\ref{11ex17}) taken from an online forum, than for the D-forms\is{D-form} to enter formal register.

As the official U-forms ``trickle down'' into less formal registers and contexts, this leads to partial semantic bleaching of the honorific \emph{pani}. It may be retained in contexts where little or no marking of respect is needed or intended. The appositional\is{apposition} structure with the word is simply the most certain way to assure clarity of the \isi{feminine} reference. As a result, this polite word becomes an auxiliary noun, an analytic marker of the \isi{feminine} \isi{gender} and carrier of the inflectional\is{inflection} marking, cf. the discussion of the data from Table \ref{11tab4}.

(\ref{11ex32}) mentions a school principal, first giving the \isi{masculine} noun (strictly speaking an \isi{M-form}) and then correcting that the principal is actually a woman (with a U-form in \isi{apposition}). The word \emph{pani} is used to mark the \isi{gender} rather than show respect. A plain U-form would feel inadequate, in line with the belief of the speakers that U-forms are ``\isi{masculine}'', cf. Section \ref{11sec3.1}.

\begin{exe}
	\ex \label{11ex32}
	\gll Mamy spoko dyrektor-a tzn. pani-ą dyrektor. \\
	have.\textsc{1pl} ok principal-\textsc{acc.sg} i.e. Mrs-\textsc{acc.sg} principal[\textsc{acc.sg}] \\
	\trans ‘We have a good principal, that is Ms principal.’ \hfill (NKJP)
\end{exe}

The same effect can be observed in metalinguistic comments. The title of the article by \citet{szpyra-kozlowska2019} quoted above presents various options for labelling a female prime minister: a \isi{D-form} created by conversion (\emph{premier-a}), a \isi{D-form} created by suffixation (\emph{premier-k-a}) and a U-form. The last one is tellingly presented in \isi{apposition} with the honorific.

Similarly, (\ref{11ex12}) above is a newspaper headline. The use of the \isi{D-form} \emph{ministra} was criticised in the comments on the newspaper's website. One of them provided the ``nicer'' alternative: \emph{pani minister}. The actual headline with the \isi{D-form} does not need a honorific and neither would a similar headline about a man.

\subsubsection{Ideology} \label{11sec3.3.5}

{\sloppy In recent years the competition of U-forms and D-forms\is{D-form} has gained ideological overtones in the public discourse in Poland. \citeauthor{stepien2018} (\citeyear{stepien2018}, after \citealt[158]{szpyra-kozlowska2021}) shows that D-forms\is{D-form} are more common in liberal and left-wing media than in right-wing ones. Particularly well publicised instances of use of individual D-forms\is{D-form} spark heated debates, cf. Section \ref{11sec4.4}. The divide can be seen in the Sejm as well. On several occasions female MPs requested their plaques and official items like voting cards bear the \isi{D-form} label of their position: \emph{posłanka} rather than U-form/\isi{M-form} \emph{poseł}.\footnote{\url{https://oko.press/kancelaria-sejmu-nie-zmienila-tabliczek-poslanki-sa-poslami} \mbox{Accessed 27.06.2024.}\label{11ftn26}} Only left-wing MPs made such requests.\par}

Word type choices can lead to arguments. Let us recall one such event from the Sejm. It is customary for MPs to start their speeches by addressing the person presiding over the session. In 2019, the MP Krzysztof Mieszkowski\ia{Mieszkowski, Krzysztof} started his speech with the phrase \emph{Pani Marszałkini}, thus using a \isi{D-form} of \emph{marszałek} `speaker of the parliament'. The presiding deputy speaker Małgorzata Gosiewska\ia{Gosiewska, Małgorzata} immediately protested with passion and requested that a U-form be used when she is addressed.\footnote{\url{https://wiadomosci.onet.pl/kraj/spiecie-w-sejmie-marszalek-gosiewska-ja-sobie-nie-zycze/thz3bdx} Accessed 23.06.2024.} The said MP came from the then liberal opposition and the deputy speaker is a member of the then ruling conservative party.\is{U-form|)}

\section{Historical development of U-forms\is{U-form}} \label{11sec4}

\subsection{Background} \label{11sec4.1}

In earlier stages of Polish a novel usage of a concept with reference to a woman rather than a man resulted in feminative derivation. Thus, D-forms\is{D-form} dominated, both among common nouns and surnames. Interesting proofs of the strength of the pattern are the nouns\emph{ papieżyca} `popess' (limited range of application), \emph{szpiegini} `spy, \textsc{f}' or \emph{gościni} `guest, \textsc{f}' (both later supplanted by M-forms\is{M-form}, cf. Section \ref{11sec3.3.3}). However, M-forms\is{M-form} (U-forms?\is{U-form})\footnote{In the \isi{case} of historical, and esp. metalinguistic evidence, it is often difficult to distinguish these two.} were used too, not necessarily very seldom. Apparently, the oldest attested \isi{M-form} in a text written in Polish comes from Mikołaj Rej\ia{Rej, Mikołaj}, a 16\textsuperscript{th} century writer \citep[65]{nitsch1951}.

Below we sketch the history of the development of U-forms\is{U-form} and their competition with other types, but let us recall that changes affected only a relatively small part of the lexicon. Numerous common nouns have been D-forms\is{D-form} for centuries, e.g. most of those in (\ref{11ex27}). Women's surnames in \emph{-sk(a)} have also been stable as D-forms\is{D-form}. Such words are hardly mentioned in the following sections. Yet, the reader should bear in mind that they were present in the system and co-existed with the changing forms we focus on.

\subsection{Surnames} \label{11sec4.2}

As for surnames, D-forms\is{D-form} were originally the norm and involved a differentiation between married and unmarried women, cf. Table \ref{11tab3}. Women were identified by the relationship to their fathers or husbands and the possessive flavour to these forms may be felt by some speakers even today. The \isi{masculine} variant was considered basic, a surname being a reflection of patrilineal ancestry. Hence, occasional mentions of women with this ``true'' form can also be found. These were probably U-forms\is{U-form} from the very beginning.

\citet[63]{kucala1978} lists examples of female surnames from a 17\textsuperscript{th} century diary, most of them are D-forms\is{D-form}, but there are several U-forms\is{U-form} as well. The practice of simplifying women's surnames into U-forms\is{U-form} gradually gained ground and allegedly was fashionable in the 18\textsuperscript{th} century Warsaw \citep[63]{nitsch1951}. By the mid 19\textsuperscript{th} century it became the \isi{subject} of heated criticism by language purists\is{language purism}, which is a proof of its spread. The criticism focused mostly on young women, esp. adolescent schoolgirls, who allegedly dared to use U-forms among themselves and with their female teachers. However, very similar complaints are attested from 1850s until 1900s \citep[50]{pawlowski1951}, which shows that U-forms\is{U-form} and D-forms\is{D-form} of surnames co-existed for decades in actual usage, most likely with register differences.

Another source of support for U-forms\is{U-form} was foreign administration of partitioned Poland, which operated in \ili{German} or \ili{Russian} and did not use traditional Polish D-forms\is{D-form} for surnames. The Austrian partition adopted Polish as the official language in the 1860s, but the use in documents in the following decades varied \citep{kowalska2013}. The region seems to have clung to \isi{D-form} surnames more than the other two partitions (\citealt[63]{nitsch1951}; \citealt[56]{pawlowski1951}).

After the First World War, the newly independent Poland introduced mandatory use of traditional, ``patriotic'' \isi{D-form} versions of female surnames in documents \citep{dacewicz2021}. Actual practice varied and D-forms\is{D-form} proved difficult in use. With some surnames they involved sound alternations, which could lead to mistakes in recreating the male versions from the D-forms\is{D-form}. Consequently, in 1930 the Ministry for Education relaxed the regulations -- in school documentation D-forms\is{D-form} remained mandatory for surnames in \emph{-wicz} \citep[54]{pawlowski1951}, otherwise U-forms\is{U-form} were allowed. Let us also recall that with most surnames, D-forms\is{D-form} revealed the marital status, which was certainly felt as a flaw of this structure by many women.

After the Second World War, U-forms\is{U-form} were once again introduced as the default versions of women's surnames in administration and official documents. Usage varied and the new standard was more eagerly used by younger people. In the early 1950s D-forms\is{D-form} did not seem under threat in informal use \citep[64]{nitsch1951}, while actively used U-forms\is{U-form} were typical of schools \citep[49]{pawlowski1951}. Pawłowski presents surveys conducted in secondary schools and among bank clerks. The data suggest differences in use based on \isi{gender} and age. Young women (pupils, bank clerks, also teachers) generally preferred U-forms\is{U-form}, while male pupils opted for traditional D-forms\is{D-form} \citep[54--55]{pawlowski1951}. Interestingly, \citet[49]{pawlowski1951} suggests that women's surnames could have the form identical to male versions in the \isi{nominative} (i.e. in the documents), but inflect\is{inflection} as D-forms\is{D-form} otherwise. It is unclear whether such hybrid usage was common at the time. It may well be just a compromise proposal by the scholar disapproving of \isi{uninflectedness}, but aware of its spread.

With time, U-forms\is{U-form} gradually dominated. \citet{kucala1978} recounts several surveys of surnames in the press, which clearly show that D-forms\is{D-form} were used less and less. In the mid 1970s, traditional D-forms\is{D-form} constituted 24\% of women's surnames in press obituaries, 6\% in a women's magazine and were absent from a youth magazine \citep[68]{kucala1978} reflecting the steady decline across generations. The prevalence of D-forms\is{D-form} in speech was probably higher, but the pattern was clearly losing ground. This tendency has never been reversed\footnote{We should note a trend to use D-forms\is{D-form} of surnames among women in artistic (maiden name versions) and scholarly (married versions) circles. The names of some authors cited here: Grzegorczykowa, Puzynina and Waszakowa represent the latter group. The fashion is now gone -- the youngest followers of it were born in the 1960s.} and today D-forms\is{D-form} for surnames are mostly spoken and informal. Yet, even within this limited range, the suffix \emph{-ow(a)} remains productive, cf. the informal \emph{Clintonowa} `Mrs. Clinton' or \emph{Merkelowa} `Mrs. Merkel'. When using such forms many speakers today apply the suffix \emph{-ow(a)} regardless of the marital status of the women involved.\footnote{This practice is parallel to the standard usage of surnames in \ili{Czech}, where most surnames have a single \isi{D-form} in \emph{-ová} (or \emph{-ská}).}

\subsection{Common nouns} \label{11sec4.3}

The ``original'' situation presented in Section \ref{11sec4.1} seems to have been the \isi{case} for common nouns for a long time. D-forms\is{D-form} dominated, appearing in the vocabulary when needed, but there were instances of M-forms\is{M-form}, esp. in the copular construction (the earliest known \isi{M-form} in Rej's\ia{Rej, Mikołaj} writings, cf. Section \ref{11sec4.1}, appears in such a context). Most likely there was some form of ``division of labour'' between the types (cf. Section \ref{11sec3.3.3}) and a minority of personal nouns were typically M-forms\is{M-form}.

The late 19\textsuperscript{th} century saw an unprecented rise of women professionals as various occupations, positions and study programmes at universities were gradually opened to them. These social changes continued in the 20\textsuperscript{th} century. The language needed labels for women doctors, lawyers, officials, etc. and these were often called by means of U-forms\is{U-form}/M-forms\is{M-form}, which spurred criticism by language purists\is{language purism} roughly since the 1880s \citep[28]{jezykpolski1933}. The commentaries were emotional and used harsh language, e.g. accusing women with professional titles of cowardly and fraudulently hiding their femininity \citep[118]{poradnikjezykowy1911}, cf. also fn. \ref{11ftn32}.

Such metalinguistic criticisms and later linguistic papers (notably \citealt{klemensiewicz1957}) give the impression that the progressive female professionals rejected D-forms\is{D-form}, while conservative purists\is{language purism} (mostly men) advocated this traditional structure. As the use of U-forms\is{U-form}/M-forms\is{M-form} was more and more common over decades, the protests allegedly became weaker as these structures gained some level of reluctant acceptance still in the interwar period \citep[112--113]{klemensiewicz1957}.

This narrative seems to present a partial picture only. First of all, it is not a given that the criticisms by individual language purists\is{language purism} and conservatives, however vocal, are in any way representative of broader social perception of the structures, let alone the usage at the time. We should also note that some language authorities were tolerant towards M-forms\is{M-form} \citep[30--31]{jezykpolski1931-188}. We have surviving metalinguistic evidence from the critics of M-forms\is{M-form}/U-forms\is{U-form}, but very little from the other side. \citet{wozniak2014} cites copiously from feminist\is{feminism} press of the interwar period, showing that D-forms\is{D-form} and M-forms\is{M-form} (she quotes no U-forms\is{U-form}) were both used. In fact, it was normal to see a mix of the types in a single text. Individual women may have had their preferences, but the women's movement as a whole did not advocate any of the word types. \citet[103]{klemensiewicz1957} recounts a 1929 census of teachers in which D-forms\is{D-form} and U-forms\is{U-form}/M-forms\is{M-form} for various positions in the educational system appear in comparable numbers\is{number}. This suggests that differences in register between the types were relatively slight. We know of no official policy of the Polish state with reference to labelling women professionals.

An alternative scenario is that the early decades of the 20\textsuperscript{th} century were a time of a constant rise in the \isi{number} of women to whom the new words could realistically apply. It is then likely that this meant that both D-forms\is{D-form} and U-forms\is{U-form}/M-forms\is{M-form} were used more and more often, but only the latter growth was noticed and discussed at length. Having been in use for several decades, the new formations naturally gained some acceptance. Furthermore, even though D-forms\is{D-form} had been well established as a pattern, each new word of the type needed to be accepted by the speakers, esp. because of the variety of feminative suffixes (\citealt[298--300]{wozniak2014}; \citealt[1]{obrebska-jablonska1949}).

\largerpage
The concept that was at the forefront of the changes in the language and the society was the figure of a woman doctor. Female medics were an increasingly common sight, had contact with many people and the position had the prestige earlier reserved for men. Apparently, most of them preferred the \isi{U-form}/\isi{M-form} \emph{doktor}, the argument being that this is the right translation of the \ili{Latin} title they had been given \citep[118]{poradnikjezykowy1911}. The purist\is{language purism} critics disapproved of the plain \isi{U-form},\footnote{The editors of \citealt{poradnikjezykowy1925-zagadnienia}: 15) described a sign advertising a medical practice of a dermatologist \emph{dr I Krzywy}. The acronym stands for the word \emph{doktor} and the surname is not a \isi{D-form} so it was perfectly possible to conclude that the said medic is a man, which was not the \isi{case}. The comment criticises the sign in a dramatic tone: what a misunderstanding it would be if a male patient decided to consult Dr. Krzywy!\label{11ftn32}} and even of the \isi{apposition} with \emph{pani}, today standard, well into the 1930s \citep[159]{jezykpolski1931-208}.\footnote{Recounting his own stay in hospital, \citet[184]{benni1933} suggested a hybrid formation, the \isi{feminine} honorific and an \isi{M-form}: \emph{z pani-ą doktor-em} `with [Ms.] doctor'. \citet[37]{klemensiewicz1957} and \citet[65]{nitsch1951} present this structure as one in use in the interwar period and even later (cf. also \citealt[374]{kupiszewski1967}), but Benni's own mention	of the surprise of the nurses and a response from \citet{obrebska1933} suggest it was not particularly frequent.} Even though the vocal purists\is{language purism} advocated the \isi{D-form} alternative \emph{doktorka}, which was in use, not all people concerned with language correctness agreed, thinking the form is either ``depreciative'' \citep[55]{poradnikjezykowy1906} or ``not strictly grammatical'' \citep[104]{poradnikjezykowy1925-odpowiedzi}.

We may conclude that U-forms\is{U-form}/M-forms\is{M-form} were becoming more common in the interwar period, but they did not gain acceptance of refined and traditional speakers. Alternative D-forms\is{D-form} were not accepted straightforwardly either. The data from \citet{klemensiewicz1957} and \citet{wozniak2014} quoted above suggest the register variation between the types was not very strong at the time. Yet, the example of female doctors suggests the modern opposition of formal/polite \isi{U-form} with the honorific vs. the informal \isi{D-form}, as described in Section \ref{11sec3.3.4}, was forming already.

The communist\is{communism} rule installed in Poland in 1945 brought important changes to society, including massive employment of women. Consequently, there were even more real-life situations when they could be referred to by their professions or addressed by professional titles. The field of competition between D-forms\is{D-form} and U-forms\is{U-form}/M-forms\is{M-form} grew immensely. The official ideological position that men and women are equal translated into the use of exactly the same labels for both sexes in official language. In other words, M-forms\is{M-form}/U-forms\is{U-form} were promoted against ``backward'' D-forms\is{D-form}. Plain U-forms\is{U-form} started to appear in written texts other than official records \citep[3--4]{obrebska-jablonska1949} and gained momentum. The move away from D-forms\is{D-form} might have been modelled on \ili{Russian} \citep[305]{wozniak2014}, but some see the development in Polish as independent \citep[250]{lazinski2006}. The use of different formations for common nouns referring to women varied in the early years after the war and so did the views of the speakers. We have surviving mentions of disagreements between speakers (e.g. \citealt{doroszewski1951}) and criticisms of \isi{uninflectedness}.

Varying opinions were expressed by linguists of the time. \citet[60, 62]{pawlowski1951} and \citet[65]{nitsch1951} accept the lack of inflection\is{uninflectedness} only in the appositional\is{apposition} phrase \emph{pani doktor}, condemning similar use of other nouns and plain U-forms\is{U-form}. Yet, we must conclude that such usage (today standard) was clearly present in the language if it was criticised. Taking an opposite ``progressive'' stance, \citet{klemensiewicz1957} announces U-forms\is{U-form}/M-forms\is{M-form} to be finally victorious in their long competition against D-forms\is{D-form} among words for (prestigious) professions and professional titles. As we shall see below, his conclusions were definitely too hasty. The supposed triumph was also limited, as a considerable portion of the \isi{feminine} lexicon had always been served by D-forms\is{D-form}, cf. Section \ref{11sec3.3.3}, and this remained so. A decade later, \citet{kupiszewski1967} also sees the use of M-forms\is{M-form}/U-forms\is{U-form} as a strong and unstoppable tendency. Interestingly, his paper is inspired by a letter from a radio listener who protested againt the trend and insisted on using the ``correct'' D-forms\is{D-form}.

\citeauthor{kucala1978}'s (\citeyear[51--61]{kucala1978}) survey of noun usage in press obituaries and announcements provides numerous examples of both D-forms\is{D-form} and M-forms\is{M-form} (more often than U-forms\is{U-form}). \citet[61]{kucala1978} notes also the use of plain U-forms\is{U-form}, but considers them to be colloquial. By the 1970s the use of U-forms\is{U-form} was accepted by prescriptivist works -- at least for those words where D-forms\is{D-form} had not become established \citep[111]{buttlerkurkowskasatkiewicz1973}, although care was advised to avoid plain U-forms\is{U-form} \citep[116]{buttlerkurkowskasatkiewicz1973}. Later prescriptivist works are still slightly more approving of U-forms\is{U-form} \citep[128--129]{jadacka2006}.

\subsection{Recent developments} \label{11sec4.4}

{\sloppy After the overthrow of \isi{communism} in 1989, the incipient feminist\is{feminism} movement started advocating for D-forms\is{D-form} instead of U-forms\is{U-form}/M-forms\is{M-form} among common nouns. For some time such usage was limited and perceived to a large extent as an ideological declaration \citep[278]{lazinski2006}. Gradually the ``new'' D-forms\is{D-form} started to appear in public discourse and raise controversy among native speakers, by now used to the presence of U-forms\is{U-form}. In recent years there were several occasions when the use of specific D-forms\is{D-form} in public statements drew great attention and led to heated debates among native speakers. The particularly controversial words were e.g. \emph{ministra} (vs. \textsc{m} \emph{minister}), \emph{powstanka} (vs. \textsc{m} \emph{powstaniec} `insurgent'), \emph{gościni} (vs. \textsc{m} \emph{gość} `guest') and \emph{posłanka} (vs. \textsc{m} \emph{poseł} `MP', cf. Section \ref{11sec3.3.5}).\par}

We will focus on the controversy with the first word. In the early 2012, the then Minister for Sport, Joanna Mucha\ia{Mucha, Joanna}, declared in a television interview that she prefers to be addressed \emph{Pani Ministro}, i.e. with a \isi{vocative} of the \isi{D-form} \emph{ministra}, as in (\ref{11ex12}). This provoked a flood of criticism, even from fellow ministers, but there were also voices of approval. Within days after the interview, a well-known linguist talked of Ms. Mucha\ia{Mucha, Joanna} ``raping the language''.\footnote{\url{https://www.wprost.pl/kraj/308115/ministro-mucho-nie-mozna-gwalcic-jezyka.html} \newline Accessed 28.06.2024\label{11ftn33}} There are also at least two linguistic articles with Ms. Mucha's\ia{Mucha, Joanna} name in the title \citep{bobrowski2012, waszakowa2014}, much more moderate in tone. Most critics disapproved of any \isi{D-form}, others would have accepted one created by suffixation: \emph{ministerka}, cf.~the analysis of the title of \citet{szpyra-kozlowska2019} in Section \ref{11sec3.3.4}. The Council for the Polish Language, an authoritative body on matters of language correctness, felt compelled to issue a statement. Their position \citep{rjp2012} was moderate, acknowledging the possibility of conversions like \emph{ministra} but voicing scepticism about their chances of being generally accepted.\footnote{Interestingly, the verdict was perceived as a victory for the \isi{D-form}, since it was not fully condemned, cf. a title of a press report stating simply that ``\emph{ministra} is correct'' \url{https://www.newsweek.pl/polska/ministra-w-jezyku-polskim-jest-poprawna-rada-jezyka-polskiego/9b9q25g}	Accessed 28.06.2024.} They also mentioned numerous messages they had received and asserted that most women were sympathetic to Ms. Mucha's\ia{Mucha, Joanna} \isi{D-form}. One of the largest Polish dailies, \emph{Gazeta Wyborcza}, uses the word \emph{ministra}, cf. (\ref{11ex12}). It seems this practice started after Ms. Mucha's\ia{Mucha, Joanna} interview.

Despite fervent opposition from purists\is{language purism} and conservatives,\footnote{Ironically, at least until the 1960s a conservative approach was to use and support D-forms\is{D-form} rather than U-forms\is{U-form}, cf. Section \ref{11sec4.3}.} D-forms\is{D-form} keep spreading in the public sphere. Occasionally, D-forms\is{D-form} appear in official contexts as well, cf. the title of the dean of one of the faculties at Adam Mickiewicz University in Poznań: \emph{dziekana} (vs. standard \textsc{m} \emph{dziekan}),\footnote{\url{https://anglistyka.amu.edu.pl/strona-glowna/wladze-wydzialu} Accessed 28.06.2024.} or the title of the deputy mayor of Sopot: \emph{wiceprezydentka} (vs. \textsc{m} \emph{wiceprezydent}).\footnote{\url{https://trojmiasto.wyborcza.pl/trojmiasto/7,35612,28938409,magdalena-czarzynska-jachim-zastepczyni-prezydenta-sopotu.html} Accessed 28.06.2024.}

As can be seen, the pace of change is quick and the more often D-forms\is{D-form} are used, the more acceptance they eventually gain. Paulina Mikuła\ia{Mikuła, Paulina}, a youtuber specialising in language and language correctness, made two videos concerning D-forms\is{D-form}. In the later one (made in 2019), she is more approving and admits that she has changed her attitude, she ``simply got used to'' D-forms\is{D-form}.\footnote{\url{https://www.youtube.com/watch?v=TjJZ8THmtXQ} Accessed 28.06.2024.} The earlier video was created only 5 years prior.\footnote{\url{https://www.youtube.com/watch?v=VJZm2vAh_Oc} Accessed 28.06.2024.}

We must note that the current debate about D-forms\is{D-form} completely ignores surnames, here the use of U-forms\is{U-form} is not contested and no change in usage is detectable. Interestingly, both the dean from Poznań and the deputy mayor from Sopot mentioned above have surnames that are U-forms\is{U-form} according to the current norm, \emph{Pawelczyk} and \emph{Jachim} respectively, and it is these forms that appear officially next to the innovative D-forms\is{D-form} for their positions.

\is{U-form|(}
\subsection{U-forms: the history of a category} \label{11sec4.5}

U-forms appeared in Polish as an innovative pattern in those cases\is{case} where the general tendency to create D-forms\is{D-form} for \isi{feminine} referents was overcome by the need to preserve the ``true'' form of the words in question, i.e. that identical to the \isi{masculine} version. The resultant conflict of morphonology and semantics made most such formations uninflected\is{uninflectedness}, but not all: M-forms\is{M-form} are attested today and were very likely more common than U-forms at early stages of the diversion from D-forms\is{D-form}. As Sections \ref{11sec4.2} and \ref{11sec4.3} show, the development happened twice, first for some types of surnames (in the period from the 18\textsuperscript{th} century until the second half of the 20\textsuperscript{th} century) and then for common nouns denoting prestigious professions, positions and professional titles (in the period from the late 19\textsuperscript{th} century until the second half of the 20\textsuperscript{th} century). In the latter \isi{case} an intermediate stage of M-forms\is{M-form} is detectable in the first half of the 20\textsuperscript{th} century.

The category of U-forms expanded in three ways: among speakers, in the lexicon and across linguistic contexts. The gradual spread in the population is presented in Sections \ref{11sec4.2} and \ref{11sec4.3}. As for the lexical range of U-forms, it has always been limited, D-forms\is{D-form} being the majority among lexical concepts\is{lexeme} with possible \isi{feminine} reference. However, U-forms successfully dominated their niches in the lexicon and were eventually accepted by the speakers for these types of words (this happened roughly in the 1960s and 1970s), to the point that the newly (re)created D-forms\is{D-form} are rejected by more conservative speakers at present. The acceptance for U-forms is still weak beyond their typical domains, cf. \citeauthor{szpyra-kozlowska2021}'s (\citeyear[336--341]{szpyra-kozlowska2021}) criticism of (\ref{11ex30}--\ref{11ex31}). Yet, attestations of such spread in the lexicon are a proof of the strength and continuing productivity of the pattern. As for the syntactic contexts, the appositional\is{apposition} construction has always been a stronghold of U-forms. It is very likely that it was their original context (with given names for surnames and the honorific \emph{pani} for common nouns). The construction remains a typical environment for U-forms, cf. Table \ref{11tab4}.

Today U-forms are used freely in the singular (except for the copular construction), but are still limited in the plural. It is likely that the hierarchy of contexts for the use in the plural shown in Table \ref{11tab5} is also a historical development path of U-forms in general: from the ``quoted'' basic form in \isi{apposition} to the fully independent use where the lack of inflection\is{uninflectedness} may be the sole marker of the \isi{feminine} \isi{gender}.\footnote{In the early 1950s, the use of the U-form \emph{doktor} was accepted by conservative linguists \citep[65]{pawlowski1951} in \isi{apposition} with the honorific, the first context in Łaziński's hierarchy, but the use of a U-form with a \isi{feminine} possessive, the next one, was novel and contestable \citep[32]{doroszewski1951}.} This perspective allows us to look again at the status of U-forms in the lexicon. We may reconcile \citeauthor{lazinski2006}'s (\citeyear[227]{lazinski2006}) argumentation against U-forms as separate lexemes\is{lexeme} and the counterexamples, cf. Section \ref{11sec3.3.2}. The latter represent a more mature stage of the development which has not been reached (fully) by too many speakers yet. In the wake of the recent developments and the rebound of D-forms\is{D-form}, U-forms may well retreat and never fully mature.

It is rather sure that U-form surnames and common nouns are facing different fates in the nearest future. The former are not threatened in any way while the latter might lose territory to D-forms\is{D-form}. This is due to different attitudes and choices of the speakers regarding the two groups.

As for common nouns, we cannot tell whether the current increased use of D-forms\is{D-form} will eventually overcome the practice of using U-forms or prove a temporary fashion. Most likely the coexistence of the types will remain for long, but it might be demarcated differently. To oust U-forms from their current contexts of use, D-forms\is{D-form} would need to lose their ideological flavour so that they feel neutral (enough) for the majority
of speakers.\is{U-form|)}

\section{Conclusions}

The paper has presented how a productive uninflected\is{uninflectedness} pattern functions in Polish and how it arose in spite of the omnipresent \isi{inflection}. Current usage shows symptoms of the difficulty posed by uninflected\is{uninflectedness} forms and repair mechanisms, chiefly \isi{apposition}. The \isi{case} of U-forms\is{U-form} provides an interesting insight about the concept of systemic pressure. U-forms\is{U-form} are exceptional in Polish grammar, but they were able to emerge as a class, spread and establish their position. The pressure of the system appears to have lain dormant for several decades as they became more and more common. The current (partial) retreat of U-forms\is{U-form} is linked first of all to the awareness of speakers and their conscious choices.

A notable side effect of the existence of U-forms\is{U-form} is the new function of the \isi{feminine} honorific \emph{pani}, which is regularly used in contexts where marking respect is not needed or intended. The word is then an auxiliary noun, carrying the \isi{inflection} lacking on the \isi{U-form} and marking the \isi{gender} explicitly.

Our analysis also shows that Polish is rather tolerant to the coexistence of different patterns and subclasses of words. The situation established by the late 20\textsuperscript{th}~century involves standard use of U-forms\is{U-form} for some surnames and some common words, while surnames in \emph{-ska} and most common words referring to women remain D-forms\is{D-form}. This kind of tolerance was also a necessary part of the emergence of U-forms\is{U-form}. The above is linked to another obsevation: Polish easily allows productivity of minor inflectional\is{inflection} patterns, which we have noted several times in Section \ref{11sec2}. U-forms\is{U-form} are only one example of this.

Another feature of Polish that surfaces in the analysis is the notion of the ``basic form'' of a noun. This status of the \isi{nominative} singular is more than a convention applied by linguists and lexicographers; it seems to match the perception the speakers have. The desire to keep this ``true'' form of a word across all contexts (which may be triggered by a \isi{number} of factors, cf. Section \ref{11sec2.2}) is \isi{instrumental} for \isi{uninflectedness}, including the \isi{case} of U-forms\is{U-form}.

We may note the role of individual words in paving the way for whole categories in historical changes. It seems that the first successful and fully accepted \isi{U-form} was the word \emph{doktor}, cf. Section \ref{11sec4.3}. Similarly, publicised controversies concerning few words have played a significant role in bringing D-forms\is{D-form} back in recent years, cf. Section \ref{11sec4.4}. We can also assume that in the final decades of the 19\textsuperscript{th} century, the presence of \isi{U-form} surnames in the language helped the rise of uninflected\is{uninflectedness} common words referring to women.\il{Polish|)}
 
 
\section*{Abbreviations}

Less typical abbreviations used in the paper are:
\vspace*{-0.2cm}

\begin{tabularx}{.45\textwidth}{lQ}
\textsc{impers} & impersonal \\
\textsc{inan} & inanimate \\
\textsc{nnom} & non-nominative \\
\end{tabularx}
\begin{tabularx}{.45\textwidth}{lQ}
	& \\
\textsc{nvir} & non-virile \\
\textsc{prep} & preposition \\
\textsc{vir} & virile (masculine personal) \\
\end{tabularx}

%\section*{Acknowledgements}
%\citet{Nordhoff2018} is useful for compiling bibliographies.

%\section*{Contributions}
%John Doe contributed to conceptualization, methodology, and validation.
%Jane Doe contributed to writing of the original draft, review, and editing.


{\sloppy\printbibliography[heading=subbibliography,notkeyword=this]}
\end{document}
